\documentclass[11pt]{article}
\usepackage[T1]{fontenc}
\usepackage{lmodern}
\usepackage[margin=1in]{geometry}
\usepackage{amsmath,amssymb,amsthm,mathtools}
\usepackage{booktabs,tabularx,array}
\usepackage{microtype}
\usepackage{xcolor}
\usepackage{xurl}
\usepackage[colorlinks=true,allcolors=blue!45!black]{hyperref}
\usepackage{fancyhdr}
\hypersetup{pdftitle={A local moment refinement near the quartic zero-free boundary},pdfauthor={Zhongpai Gao},pdfsubject={Research draft: proposed source-level analytic proof under review}}
\pagestyle{fancy}
\fancyhf{}
\fancyhead[L]{\small Local moment refinement}
\fancyhead[R]{\small Research draft \textbullet\ 11 October 2026}
\fancyfoot[C]{\thepage}
\setlength{\headheight}{14pt}
\setlength{\emergencystretch}{2em}
\setlength{\parskip}{3pt}
\allowdisplaybreaks[1]
\newtheorem{proposition}{Proposition}[section]
\newtheorem{lemma}[proposition]{Lemma}
\newtheorem{theorem}[proposition]{Theorem}
\theoremstyle{definition}
\newtheorem{hypothesis}[proposition]{Hypothesis}
\newtheorem{remark}[proposition]{Remark}
\newcommand{\Q}{\mathbf Q}
\newcommand{\Z}{\mathbf Z}
\newcommand{\C}{\mathbf C}
\newcommand{\OO}{\mathcal O}
\newcommand{\eps}{\varepsilon}
\newcommand{\pos}[1]{(#1)_{+}}
\newcommand{\code}[1]{\texttt{\small #1}}
\title{A local moment refinement near the\protect\\quartic zero-free boundary\\[6pt]
\large A proposed analytic proof with an exact endpoint certificate}
\author{Zhongpai Gao\thanks{Prepared with AI assistance in research, computation, formalization, and exposition. The upstream mathematical contributions are credited in Section~\ref{sec:credits}. This draft has not received independent mathematical certification.}}
\date{11 October 2026 \quad---\quad Draft 3}

\begin{document}
\maketitle
\begin{abstract}
We propose a local refinement of the moment estimate in Collet's quartic
zero-free argument. The coefficient of the total prime-slot length is
reduced from $9/2$ to $2249/500=4.498$ in an explicit neighborhood of the
relevant pair of ordinary summation lengths. Rows with a smaller buffered
zero parameter are paid by removing slots and using the existing moment.
Rows with a larger parameter are counted together with their arithmetic
labels through one effective integer row. This draft supplies a proposed
source-level analytic proof: exact signed characters and masks, the fixed
character dictionary, the unmarked density input, restricted child norms,
and the common Fourier transport are assembled in the supplements.
An exact rational certificate supports the boundary
$4374433/5000000=0.8748866$. The final supplement gives a complete proposed
written assembly of the weighted numerator and continuation argument for
finite-order Eisenstein Hecke characters, ordinary Dirichlet characters,
and the Riemann zeta function. This proof is under review and relies on
the pinned incumbent analytic framework. It is not an independently
verified record. Compiled Lean results support intermediate steps;
they are not a complete formal proof of the proposed endpoint.
\end{abstract}
\clearpage
\tableofcontents
\clearpage

\section{Introduction and the status of the result}

A fixed zero-free half-plane is a statement that an $L$-function has no
zeros to the right of a specified vertical line. For the families under
discussion, lowering that line is an improvement. The eventual aim of
this project is a strict half-plane $\Re s>\theta$ for ordinary Dirichlet
$L$-functions, including the Riemann zeta function, and for finite-order
Hecke $L$-functions of trivial infinite type over $\Q(\sqrt{-3})$.
Principal poles are excluded. A bound near $0.875$ is a fixed-strip
improvement; it does not approach a proof of the Riemann hypothesis by
itself.

Our starting point is the pinned Cycle25 manuscript of Collet
\cite{collet}, which builds on OpenAI's $7/8$ framework \cite{openai}
and Levy's weighted numerator \cite{levy}. Its exact comparison boundary is
\begin{equation}\label{eq:quartic}
 b_0=\frac{11}{12}-\frac{\ell_0}{4}
     =0.8748866470399641\ldots,
 \qquad
 P(\ell_0)=0,
\end{equation}
where
\[
 P(X)=927X^4-3135X^3+2433X^2+275X-100,
 \qquad \ell_0\in[1/6,1/5].
\]
Throughout this paper, ``comparison boundary'' refers to this particular
pinned source, rather than to a claim about the current leaderboard.

The proposed modification is deliberately local. In the source proof,
two ordinary character sums and a product of short prime sums share a
moment budget. Write their logarithmic lengths as $n_1,n_2,z$ and the
row budget as $M$. The existing estimate is available with
$n_1+n_2+(9/2)z\le M$. We seek to replace $9/2$ by $4.498$ only when
the lengths lie near the configuration used by the endpoint argument.
Because $z>0$, a smaller coefficient permits a slightly larger prime
product at the same row budget. The weighted numerator can then gain a
small amount.

The proposed endpoint is
\begin{equation}\label{eq:theta}
 \theta=\frac{4374433}{5000000}=0.8748866.
\end{equation}
Its distance below \eqref{eq:quartic} is about $4.704\times10^{-8}$.
The size of this gain makes exact arithmetic essential. More
importantly, a favorable scalar inequality is useful only if it belongs
to a valid estimate for the original arithmetic sums.

The status of each assertion matters. The elementary inequalities and
rational endpoint certificate are established within their algebraic
scope. The source-level analytic argument presented here is a proposed
proof under review, relative to the pinned incumbent theorems. Several of
its intermediate identities and estimates have compiled Lean proofs using
cached dependencies. A complete new formal endpoint theorem and independent
verification are absent. The all-family conclusion at
\eqref{eq:theta} is the proposed theorem in
Section~\ref{sec:endpoint-assembly}, with a complete written argument
supplied for review. Its proof status must be distinguished from
independent acceptance or formal certification.

Here is the roadmap. Section~\ref{sec:moment} states the exact family and
local estimate. Sections~\ref{sec:idea}--\ref{sec:interfaces} explain its two
branches and the reserves paying the remaining terms. Section~\ref{sec:endpoint}
gives the conditional numerical consequence. The supplements then present
the detailed source arguments in dependency order: signed reconstruction,
the fixed character dictionary, plain density, uniform Fourier transport,
restricted canonical children, the assembled local moment, and shortened
physical scales. The final supplement assembles the full numerator and
all-family continuation, including short selected products and every
free Euler shift. The independent-validation scope is stated in
Section~\ref{sec:gaps}.

\begin{center}\small
\begin{tabularx}{\textwidth}{@{}p{0.27\textwidth}X@{}}
\toprule
Input or conclusion & Role in the argument\\
\midrule
Pinned $7/8$ and $9/2$ results & Existing analytic inputs; called on children with their actual masks and moving characters.\\
Signed row and fixed dictionary & Identify the character being counted and preserve literal eligibility.\\
Low-row payment and high-row density & The two parts of the proposed local $4.498$ moment.\\
Exact rational certificate & Converts that moment, if valid in the full numerator, into a strict endpoint saving.\\
All-family endpoint & Proposed complete written proof relative to the pinned framework; independent review and complete formal verification remain outstanding.\\
\bottomrule
\end{tabularx}
\end{center}

\section{The proposed local moment}\label{sec:moment}

\subsection{Objects and normalization}
We use the natural moment family of \cite[Appendix B3.4]{collet} over
$\OO=\Z[\omega]$, where $\omega=e^{2\pi i/3}$. Write $N$ for the
ideal norm, or the norm of an element. Set $U=Z^M$, with $Z>1$ and
$M>0$. At the parent stage the character acting on a summation variable
$n$ has the form
\[
 \psi_k(n)=\tau(n)\chi_n(k),
\]
where $\tau$ is a fixed finite ray character and $\chi_n(k)$ is the
source's sextic residue character, with its prescribed zero extension.
The ray modulus is fixed independently of $U$. No additional parent
character of positive moving width is allowed in our proposed estimate.
Such factors can occur in transformed children and must still be paid.

Let $S_{k,1}(Z^{n_1})$ and $S_{k,2}(Z^{n_2})$ denote the source's
centrally normalized ordinary character sums. The two smooth profiles,
real norm powers, and height parameters may be independent. A prime
slot has the form
\begin{equation}\label{eq:slot}
 Q_{k,i}=Z^{-z_i/2}
 \sum_{\mathfrak p\in\mathcal P_i^{\rm can}}
 \psi_k(\mathfrak p)\nu_i(\mathfrak p)
 W\!\left(\frac{N\mathfrak p}{Z^{z_i}}\right)
 \left(\frac{N\mathfrak p}{Z^{z_i}}\right)^{-\sigma_i+iv_i}.
\end{equation}
Here $\mathcal P_i^{\rm can}$ is the genuine canonical good-prime pool,
$\nu_i$ is an allowed relative ray character, and the slots share the
one smooth annular profile $W$. Put $z=\sum_i z_i$ and
$Q_{k,I}=\prod_{i\in I}Q_{k,i}$. Formula~\eqref{eq:slot} does not
permit arbitrary prime subsets or independently chosen slot profiles.

An additional common puncture is allowed: all factors may be restricted
to arguments coprime to a fixed nonzero mask of norm at most $U^B$.
It stays fixed during the current row sum but may depend on labels
already frozen in a preceding transformation. The underlying natural
zeros must also stay in place. A canceled character phase at a nonunit
does not turn the original value zero into one.

The word \emph{eligible} means the literal source predicate excluding
\code{FixedInducingRow} presentations. Let $M_{\rm ray}$ be the fixed
ray modulus and $\eta_0$ the fixed anchor in the relative slot characters.
The initial source uses
\[
 Q_* = M_{\rm ray}\cap(72)\cap\mathfrak q_{\eta_0}.
\]
This ideal is held fixed in transformed states, even when their
whole-product character moves. Its containments cover every allowed
relative ray twist; they are part of the initial family, not a later
change in the exceptional set. Eligibility is not replaced by
nonprincipality of one preferred primitive representative. Its exceptional
part is treated by the centered cancellation. The exact transport is
spelled out in Section~\ref{sec:fixed-dictionary-bridge}.

\begin{proposition}[Local physical moment: proposed proof under review]\label{hyp:local}
Use the preceding family, with all its natural masks and canonical slot
conditions. Suppose
\begin{equation}\label{eq:window}
 \frac{41M}{100}\le n_i\le\frac{11M}{25},\qquad
 \frac{3M}{100}\le z\le\frac M{25},\qquad
 n_1+n_2+Cz\le M,\qquad C=\frac{2249}{500}.
\end{equation}
For every $\eps>0$, the proposed estimate is
\begin{equation}\label{eq:moment}
 \sum_{k\ {\rm eligible}}
 \left|S_{k,1}(Z^{n_1})S_{k,2}(Z^{n_2})Q_{k,I}\right|^2
 \ll Z^{M+\eps}\mathcal H_J.
\end{equation}
The source's radial row weight is understood. The factor $\mathcal H_J$
denotes its finite smooth-seminorm and polynomial height allowance,
retaining independent heights. The assertion also includes the four
annular derivative-profile pairs and bounded scale-box recovery used by
the physical numerator. All quantifiers, including selection of a common
fine mesh before the later arithmetic data, are as in the source family.
\end{proposition}
\begin{proof}[Proof reference]
The exact source-level proof is given in
Section~\ref{sec:physical-local-assembly}, after the signed reconstruction,
fixed-dictionary, plain-density, Fourier, and restricted-child lemmas.
\end{proof}

The notation imports a substantial existing analytic framework. This
draft explains a modification of that framework; it is not a replacement
construction of the ordinary $L$-functions or of the source moment family.
Sections~\ref{sec:restricted-canonical} and~\ref{sec:physical-local-assembly}
provide the proposed connection to the actual canonical state and the
complete local source assembly. They are written analytic arguments with
compiled intermediate steps, rather than a new formal endpoint theorem.

\section{Why a local improvement might close}\label{sec:idea}

There are two potential difficulties with improving a moment estimate.
First, a transformed child need not remain in the same local length
window as its parent. Second, counting a character family after freezing
its arithmetic labels may be too expensive. The proposed argument treats
these difficulties separately.

Apply the two transformations in the existing moment proof and classify
the resulting nonexceptional rows by a buffered zero parameter $\delta$.
Roughly, the associated real zero location is $(1+\delta)/2$; the
buffer and finite character family are those of the source. When
$\delta\le7/10$, pointwise prime estimates make it cheap to remove
enough slots to recover the existing global moment condition. For
$\delta\ge7/10$, those pointwise estimates are more expensive, but the
rows are sparse. The proposed saving comes from counting each such row
\emph{together with} its common-support labels.

The split is made after the transformations. No Poisson formula is
applied to a zero-bin indicator. Also, the smaller-parameter branch calls
the existing moment rather than calling Proposition~\ref{hyp:local} on a
shortened child. Thus there is no claimed local induction whose length
window must somehow survive indefinitely.

\subsection{The size of the budget defect}
Define
\begin{equation}\label{eq:budget}
 A=n_1+n_2+z,\qquad
 g=\pos{n_1+n_2+\tfrac92z-M},\qquad
 a=\frac M{2000},\qquad \sigma=3a.
\end{equation}
The auxiliary amplifier has prime norms on the scale $Z^a$, disjoint
from the live slot windows after the mesh is chosen sufficiently fine.
Its length $a$ is fixed relative to $M$; it is not an arbitrarily small
final error.

\begin{lemma}[Elementary parent bounds]\label{lem:parent}
Under \eqref{eq:window},
\[
 0\le g\le\frac M{12500},\qquad
 \frac{85M}{100}\le A\le\frac{44753M}{50000}<\frac{9M}{10}.
\]
\end{lemma}
\begin{proof}
The difference of the two capacity coefficients is
$9/2-C=1/500$. Hence $g\le z/500\le M/12500$.
The lower bound for $A$ follows by adding the lower length bounds.
For the upper bound, use
\[
 A\le M-(C-1)z
 \le M-\frac{1749}{500}\frac{3M}{100}
 =\frac{44753M}{50000}.
\]
\end{proof}

\subsection{Rows with smaller zero parameters}
This branch uses the old moment on the actual transformed state. Let
$M_{\rm act}$ be that state's width: the logarithm of its dyadic row scale
plus the logarithm of its displayed character modulus. Let $M_{\rm nom}$
be the nominal width in the exponent ledger. They need not be equal.
For actual clipped plain lengths $x_1,x_2$ and remaining slot widths $z_i$,
put
\[
 F_{\rm act}=\pos{x_1+x_2+\tfrac92\sum_i z_i-M_{\rm act}}.
\]
One greedy selection removes a total slot width at most
$F_{\rm act}/(9/2)+\eta_{\rm mesh}$ until the old capacity condition holds
or every slot has been removed. On the selected low rows, the buffered
prime estimate charges $k_0=7/10$ per unit width in squared norm.
Only after this payment is the low-row restriction dropped by positivity.
The remaining norm is exactly the canonical norm to which the old
moment applies. If no slot remains, its unrestricted zero-slot estimate
has no further plain-capacity condition.

The comparison of actual and nominal widths is essential. Set
$d=M_{\rm act}-M_{\rm nom}$, with $d\le\eta_2$, and suppose the nominal
defect is bounded by $F_{\rm nom}\ge0$. Then
$F_{\rm act}\le\pos{F_{\rm nom}-d}$. For $0\le r\le1$,
\begin{equation}\label{eq:widthpayment}
 d+r\pos{F_{\rm nom}-d}\le\eta_2+rF_{\rm nom}
 \qquad(\eta_2\ge0).
\end{equation}
This follows by considering $F_{\rm nom}-d\ge0$ and its complement.
Thus a smaller actual width offsets its larger capacity defect. One must
not discard this offset and claim that the actual defect itself is bounded
by the nominal one.

For the main low branch, $r=k_0/(9/2)=7/45$ and
$F_{\rm nom}\le g+6a+\eta$, where $\eta$ includes the enlarged-ball and
support padding. The exact source child saving is at least $a$. Therefore
\begin{equation}\label{eq:lowreserve}
 a-\frac{k_0}{9/2}(6a+g)
 \ge\frac M{2000}-\frac7{45}
       \left(\frac{3M}{1000}+\frac M{10000}\right)
 =\frac{4M}{225000}>0.
\end{equation}
Choose the mesh, buffer and aggregate support losses within this fixed
reserve. The actual coefficient factorization and radial restriction,
including why no indicator is Poisson-summed, are proved in
Section~\ref{sec:restricted-canonical}. The compiled results check the
restriction, the greedy choice, and inequality~\eqref{eq:widthpayment}.

\section{Counting rows together with their labels}\label{sec:high}

\subsection{From sixth-power-free rows to all rows}
The unmarked density exponent supplied by the existing source is
\begin{equation}\label{eq:rho}
 \rho(\delta)=1-\frac{10\delta}{3(5-2\delta)},
 \qquad \frac7{10}\le\delta\le\frac34.
\end{equation}
The adjective ``unmarked'' means that this count requires no positive
prime-slot supply. It does not mean that arbitrary moving ray characters
can be inserted. A single fixed finite character list is selected before
the live labels.

Every nonzero row has a sixth-power decomposition $j=kb^6$, with the
source's primary-generator and unit conventions. The extra factors of
$b$ leave Euler holes in the displayed character. In the positive
half-plane they preserve the zeros of its primitive inducing character.
This permits the following general counting step.

\begin{lemma}[Summing the sixth-power part]\label{lem:sixth}
Assume a uniform count $\ll Y^{\rho+\eps}$ for the relevant
sixth-power-free cores of norm at most $Y$, for every $Y\ge1$.
If $\rho\ge2/7$, the corresponding count of all rows of norm at most
$H\ge1$ is $\ll H^{\rho+\eps}$. The constant may be chosen uniformly
for $\rho$ in a fixed interval with lower endpoint $2/7$, provided the
core estimate is itself uniform.
\end{lemma}
\begin{proof}
Up to the fixed unit-sector convention, the count is at most
\[
 H^{\rho+\eps}\sum_{Nb\le H^{1/6}}(Nb)^{-6(\rho+\eps)}.
\]
The ideal series converges since $6\rho\ge12/7>1$, and a single
convergent series with exponent $12/7$ dominates the sum uniformly.
If the input estimate was initially only eventual, bounded windows must
first be paid with a uniform constant. This is a genuine requirement,
not an automatic consequence of an eventual estimate with varying data.
\end{proof}

Here $\rho(7/10)=19/54$ and $\rho(3/4)=2/7$, and $\rho$ is decreasing.
Fixed-inducing presentations stay excluded from this count; their
exceptional contribution must be treated separately.

\subsection{A simple branch that explains the saving}
Consider the second-transform branch in which the same squarefree label
$P$ occurs once on each side, with $\log_Z NP=q$. Its actual character
row is
\[
 j=P^2h',\qquad \log_Z Nj\le M-6a+o(1).
\]
At a fixed $j$, the labels $P$ satisfy $P^2\mid j$. There are only
divisor-many such labels, at an arbitrarily small power cost. Thus a
density bound on $j$ controls the joint population $(P,h')$ without
an additional factor $Z^q$.

This ordering matters. If one first freezes $P$, bounds the remaining
rows for each twist, and finally pays for all choices of $P$, one loses
the information that $P$ is already contained in $j$.

For this branch the source normalization leads to the proposed exponent
\begin{equation}\label{eq:simplehigh}
 A-M+5a+(M-6a)\rho(\delta)+\delta(A-q).
\end{equation}
Relative to the allowance $A$, it has excess
$M(\rho-1)+\delta A+a(5-6\rho)-\delta q$.
Equation~\eqref{eq:simplehigh} illustrates the joint count. The general
identity below, including its masks, supplies the arithmetic statement
needed beyond this illustrative branch.

\subsection{The effective integer row in the general branch}
The first common-support extraction introduces widths $c,d,p$, where
$p$ is now the first common radical width. Let $R$ be its active
conductor width and $E$ the mask-divisor width. These are local
bookkeeping variables; the $p$ of Section~\ref{sec:endpoint} will instead
denote a detector parameter. Set $T=(c-d)_+$.

The signs of the two separated twists are opposite. We fix the convention
of the formal source: $t_{\mathfrak p}$ is the left multiplicity minus
the right multiplicity, modulo six. The first coefficient contains the
active character on the left and its conjugate on the right. After
reciprocity, the exact row residues are therefore
\[
 e_C(\mathfrak p)=(1+t_{\mathfrak p})\bmod6,\qquad
 e_D(\mathfrak p)=(1-t_{\mathfrak p})\bmod6.
\]
To see the sign directly, write $c_n(\pi)$ for the zero-extended sextic
row and $q(\pi,n)$ for its fixed reciprocity phase. Reciprocity and
multiplicativity give
\[
 c_n(\pi)c_\pi(n)^t=q(\pi,n)^t c_n(\pi)^{1+t},\qquad
 c_n(\pi)\overline{c_\pi(n)^t}
   =q(\pi,n)^{6-t}c_n(\pi)^{7-t}.
\]
For $1\le t\le5$, these are identities including nonunit zeros.
An exponent six is the coprimality mask, not the constant one.
Section~\ref{sec:signed-reconstruction} gives the full conventions and
explains the unit relating an ideal generator to chosen prime generators.
Define
\[
 F_C=\prod_{\mathfrak p\mid r}\mathfrak p^{e_C(\mathfrak p)},\quad
 L_C=\log_Z NF_C,\quad
 R_{\rm zero}=\sum_{\substack{\mathfrak p\mid r\\e_C(\mathfrak p)=0}}
                       \log_Z N\mathfrak p.
\]
An exponent zero here records cancellation of a phase; its old puncture
is still retained.

The signed identity after the second common-support extraction gives
\begin{equation}\label{eq:integerrow}
 w=F_CeG_cV_{\rm id}h',\qquad
 \log_Z Nw\le M-R-T+L_C-6a+o(1).
\end{equation}
The factors $e,G_c,V_{\rm id}$ are the source's mask and second
extraction factors. The analytic row ball must actually be enlarged to
the smooth envelope
\begin{equation}\label{eq:envelope}
 H=K_0+(d-c)_++6a+\varpi,\qquad K_0=2A-c-d+R+E-M,
\end{equation}
before the second transformation. Here $\varpi>0$ pays the retained
frequency padding and the fixed support-doubling factor. If the retained
first dyad satisfies $K\le K_0+\delta_1$, take
$\varpi\ge\delta_1$ with a further positive support reserve. The amplified
rows have norm at most $Z^{K+6a}$, so the actual smooth ball contains them.
This is exactly the source's \code{mainCommonRadius} with $\sigma=3a$.
The added $\varpi$ only shortens the dual range in
\eqref{eq:integerrow}; its low-branch cost was reserved above.
For a fixed parent character modulus, write $M=m+q$ and use $m$ in the
formula for $K_0$; the same final total-width identity results.
In particular, the actual nominal child width is
\[
 M_{\rm nom}=M-(c-d)_+-6a-\varpi-g_2+t_2\le M-6a-\varpi.
\]

At fixed $w$, the reconstruction makes the active first primes,
the mask primes, the second common primes, and the first residual
projector label divisor-many. First common primes absent from $w$ must
still be counted; their width is
\begin{equation}\label{eq:missing}
 p-R-E+R_{\rm zero}.
\end{equation}
In particular, a canceled prime cannot simultaneously be omitted from
$w$ and treated as one of its divisors. A general arithmetic justification
of this counting step is now given in the companion reconstruction note
\cite{gaov2}. At a fixed row $w$, write $L$ for the radical of first
common primes absent from $w$. All the remaining labels are supported
on $\operatorname{rad}(w)L$. For a fixed bound on their logarithmic norms,
Rankin's inequality gives, for every $\eps>0$,
\[
 \#\{I:NI\le Z^T,\ \operatorname{rad}I\mid Q\}
 \ll_{B,T,\eps}Z^\eps\qquad(NQ\le Z^B).
\]
Indeed the count is at most
$Z^{aT}\prod_{\mathfrak p\mid Q}(1-(N\mathfrak p)^{-a})^{-1}$;
take $aT\le\eps/2$ and bound the product by a constant times
$(NQ)^{\eps/(2B)}$. The finite exceptional set of small primes contributes
only a constant. Counting the missing radical costs its norm exponent
\eqref{eq:missing}. This is a general reconstruction bound, rather than a finite sample of
valuations. The source identities supplying its factorization and mask
assumptions are given in Section~\ref{sec:signed-reconstruction}.
The second residual projector is different from the first one: it need
not divide $w$. Keep it fixed outside the row sum, retain its full signed
allocation, and pay its count together with both extraction scalars by
B3.67--B3.69. No extra puncture is imposed on the allocated quotients.

The first labels and the full coprimality projector remain inside the
positive norm until this joint count is applied. Their M\"obius signs
are retained through the exact allocation. A common Fourier measure is
chosen on the ambient parameter box before labels and bins; only its
evaluation points depend on the later arithmetic data. This order is
needed to keep independent heights and smooth-profile bounds uniform.

\subsection{What the proposed accounting gives}
Let $b_2,g_2,t_2$ denote the second common widths, with
$0\le g_2\le b_2$ and $t_2\ge0$, and put
$B_c=(3c-5d-R)_+/6$. Relative to the joint allowance
$p+A-c+R+E+B_c$, the proposed excess is
\begin{align}\label{eq:ledger}
 \mathcal E={}&M(\rho-1)+\delta A+a(5-6\rho)\notag\\
 &+(1-\rho)T-\delta c-(1+\rho)R-E-B_c
      +\rho L_C+R_{\rm zero}\notag\\
 &+g_2-t_2-(1+\delta)b_2+\eta.
\end{align}
The rational audit checks this algebraic ledger. Its analytic input is
the signed reconstruction and exact extraction normalization, followed by
the fixed-dictionary density estimate and common Fourier transport in
Sections~\ref{sec:signed-reconstruction}--\ref{sec:uniform-transport}.
Section~\ref{sec:physical-local-assembly} gives their proposed composition.

The favorable signs can be checked directly. Using
$c=\min(c,d)+T$, $\min(c,d)\ge p$, $R\le p$, and
$L_C\le5(R-R_{\rm zero})$, the second line of \eqref{eq:ledger} is at
most
\[
 (4\rho-1-\delta)p-E-B_c
 \le-\frac{79}{270}p-E-B_c.
\]
Indeed $\rho+\delta-1>0$, $5\rho-1>0$, and $4\rho-1>0$ throughout
the interval. Also $4\rho-1-\delta$ is decreasing and equals
$-79/270$ at $\delta=7/10$. The third line without $\eta$ is at most
$-\delta b_2-t_2\le0$.

Finally \eqref{eq:window} implies the weaker bound
$A/M\le1-(16/5)(3/100)=113/125$. Since
\[
 \rho'(\delta)=-\frac{50}{3(5-2\delta)^2},
\]
the function $\rho(\delta)+(113/125)\delta$ decreases on
$[7/10,3/4]$. Therefore
\begin{equation}\label{eq:highreserve}
 1-\rho(\delta)-\delta A/M\ge\frac{259}{16875},\qquad
 5-6\rho(\delta)\le\frac{23}{7}.
\end{equation}
For $a=M/2000$ and $\eta\le M/1000$, equations
\eqref{eq:ledger}--\eqref{eq:highreserve} give
\[
 \mathcal E\le
 -\frac{259M}{16875}+\frac{23M}{14000}+\frac M{1000}
 <-\frac M{100}.
\]
This leaves a substantial reserve for support and buffering losses.
The distinction between the scalar verification and its proposed analytic
source assembly is retained in Section~\ref{sec:gaps}.

\section{Other branches and uniformity}\label{sec:interfaces}

\subsection{Comparison terms and additional masks}
The centered source uses equal-product comparison rectangles with one
plain length $M/4$. For an additional common mask, the proposed argument
expands that mask only while bounding the comparison terms. After
reflecting the second natural daughter if its length exceeds $M/2$, the
two plain lengths sum to at most $3M/4+\xi$. Their existing-moment
budget is consequently at most
\[
 \frac{3M}{4}+\frac92\frac M{25}+\xi
 =\frac{93M}{100}+\xi<M
\]
for sufficiently small padding $\xi$. This pays comparison daughters
without asserting that they belong to the local window
\eqref{eq:window}. The mask remains on both rectangles of the original
centered expression.

\subsection{The amplifier errors}
For the three extracted valuation alternatives $1,6,7$, the relevant
width pairs are $(\lambda,\lambda)$, $(6\lambda,0)$, and
$(7\lambda,\lambda)$, with $\lambda\ge\sigma/4$. The old moment
pays these branches. The source's defect estimate is
\[
 F_{\rm nom}\le g+\sigma+w-w_o+\eta,
\]
its combined width and removal cost is at most
$\eta_2+F_{\rm nom}/6$ by \eqref{eq:widthpayment}, and its saving is at least
$w$. Substituting $g\le M/10000$ gives the following lower margins
before padding:
\begin{center}
\begin{tabular}{ccc}
\toprule
Valuation & $(w,w_o)$ & Saving minus base cost\\
\midrule
$1$ & $(\lambda,\lambda)$ & $13M/120000$\\
$6$ & $(6\lambda,0)$ & $193M/120000$\\
$7$ & $(7\lambda,\lambda)$ & $119M/60000$\\
\bottomrule
\end{tabular}
\end{center}
The normalized prime average must be retained. Introducing a second
count of the auxiliary primes would change these margins.

\subsection{Exceptional terms cannot be hidden in an error}
Lemma~\ref{lem:parent} gives $A-M<-M/10$. The source's second
zero-frequency term has excess at most $A-M+5\sigma/3+\eta$, and
the centered fixed-inducing calculation retains a maximum $A-M$ before
its terminal loss $3\sigma$. Both fit within this negative margin.
Replacing $A-M$ by $(A-M)_+$ would discard the saving that pays the
fixed amplifier.

For the enlarged row ball \eqref{eq:envelope}, the first exceptional
reserve is checked at that actual radius. Before the $\varpi$ padding,
the exact expressions are
\[
 F_1=\begin{cases}
 (c+5d+R+E)/6+B_c-4a,&c\ge d,\\
 c+(R+E)/6+B_c-4a,&c\le d.
 \end{cases}
\]
The padding subtracts $5\varpi/6$. Thus
$F_1\ge2c/3-4a-5\varpi/6$, while $F_2\ge2b_2/3$.
The earlier $-5a$ expression was a weaker lower bound, not an equality:
it dropped the positive amplifier-density term $a$.
The exact centered cancellation, forced congruences and raw extraction
factors are retained in Section~\ref{sec:physical-local-assembly}.
The resulting main exceptional loss $4a+O(\eta)$ fits within $M-A$;
the original error branches retain their $3\sigma+O(\eta)$ bound.

Uniformity is part of the assertion. One first fixes the requested
final loss, then chooses the aggregate small losses below both that
loss and the margin in \eqref{eq:lowreserve}. Profile orders, height
orders, and the eventual threshold are chosen through the finite chain
of transformations. A loss per slot, or a character family chosen only
after the live labels, would not supply the required estimate.

\section{The exact conditional endpoint}\label{sec:endpoint}

This section describes an independently reproducible scalar conclusion.
It uses the existing detector function; only the physical moment
coefficient is changed. Here $p\in[0,1/2]$ is a detector parameter,
unrelated to the common radical width in Section~\ref{sec:high}.

\subsection{The rational geometry and the exponent function}
Set
\begin{align}\label{eq:geometry}
 \ell&=\frac{626701}{3750000}, &
 h&=\frac{3969737987}{5000000000},\notag\\
 x&=1+\ell-h=\frac{5597590039}{15000000000}, &
 y&=h-2\ell=\frac{6895605961}{15000000000},\notag\\
 \kappa_*&=2\theta-1=\frac{1874433}{2500000}, &
 C&=\frac{2249}{500}.
\end{align}
Then $\theta=11/12-\ell/4$ and $x+y=1-\ell$.
For $\kappa\ge\kappa_*$ define
\begin{align}\label{eq:detector}
 D_\kappa(p)&=3-p-\frac{2p}{3\kappa}, &
 B_\kappa(p)&=2(1-p)\left(1-\frac{p}{3\kappa}\right),\notag\\
 c_\kappa(p)&=B_\kappa(p)/D_\kappa(p), &
 V&=\frac56-\delta,\notag\\
 R_\kappa(\delta,p)&=1-\delta+
       \frac{\delta c_\kappa(p)V}{2(V+\delta c_\kappa(p))}.
\end{align}
We use a lower-case $c_\kappa$ to distinguish the detector coefficient
from the physical coefficient $C$. The candidate upper-row exponent is
\begin{equation}\label{eq:F}
 F_{C,\kappa}(\delta,p)=\frac{1+\ell}{2}+\frac\delta2
 +h\left(\frac{3R_\kappa(\delta,p)}4-\frac5{12}\right)
 +\delta p\left(\ell-\frac{h-4\ell}{2C}\right).
\end{equation}

\begin{proposition}[Exact scalar certificate]\label{prop:certificate}
With \eqref{eq:geometry},
\[
 F_{C,\kappa_*}(\delta,p)\le\theta-10^{-7}
 \quad\text{for every }(\delta,p)\in[0,3/4]\times[0,1/2].
\]
Moreover $\theta<b_0$ for the exact boundary \eqref{eq:quartic}.
\end{proposition}
\begin{proof}[Computer-assisted proof of the first assertion]
All constants in this calculation are rational. Put
$J=VD_{\kappa_*}+\delta B_{\kappa_*}$. Then
\[
 R_{\kappa_*}=\frac{(1-\delta)J+\delta B_{\kappa_*}V/2}{J}.
\]
The stored certificate first gives $J\ge1/10$ on the whole rectangle.
After multiplying $\theta-10^{-7}-F_{C,\kappa_*}$ by $J$, it suffices
to verify a polynomial inequality. The audit partitions the rectangle
into 18 rational subrectangles and writes the polynomial in the
Bernstein basis on each one. Every coefficient is nonnegative; the
recorded closed subrectangles cover the entire original domain.
Appendix~\ref{app:reproduce} identifies the executable audit and its
exact certificate. This is an exact finite verification, not a sampling
of the exponent on a numerical grid.
\end{proof}
\begin{proof}[Proof of the comparison with the quartic]
We have $P(1/6)<0$, and exact substitution gives
\[
 P(\ell)=\frac{3484615068030068970503}
                    {21972656250000000000000000}>0.
\]
For $X\in[1/6,1/5]$,
\[
 P'(X)=3708X^3-9405X^2+4866X+275
 \ge-\frac{9405}{25}+\frac{4866}{6}+275
 =\frac{3549}{5}>0.
\]
Thus $\ell>\ell_0$, and $\theta=11/12-\ell/4<b_0$.
\end{proof}

\subsection{Why the local length restriction matters}
The row exponent in the numerator varies; it is not always $h$.
Let $d$ denote that exponent and $t$ a normalized free Euler shift.
The local estimate is intended only for
\begin{equation}\label{eq:applicationwindow}
 h-10^{-3}\le d\le h+10^{-10},\qquad 0\le t\le10^{-3}.
\end{equation}
The audit checks that both reflected plain lengths, normalized by the
row budget, lie strictly between $0.41$ and $0.44$, and that the selected
prime width lies strictly between $0.03$ and $0.04$. It also checks
that the old $9/2$ defect is less than $1/12500$ and that the original
physical slots supply the selected product.

For clarity, a general scalar row exponent used by the audit is
\begin{align}\label{eq:G}
 G_{c,\kappa}(d;\delta,p)={}&\frac12
 +h\left(\frac{17}{50}-\frac16\right)-\frac y2-\frac\ell2
 +d\left(\frac14-\frac{17}{50}\right)+\frac\delta2\notag\\
 &+\frac{3d}{4}R_\kappa(\delta,p)
 +\delta p\left(\ell-\frac{2y-d}{2c}\right).
\end{align}
It is affine in $d$ on the range in question, since $d\ge1/2>y$.
At $d=h$ and $c=C$ it reduces to \eqref{eq:F}. The exact checks are
\begin{center}
\begin{tabular}{lccc}
\toprule
Row endpoint & Coefficient $c$ & Saving below $\theta$ & Rectangles\\
\midrule
$h$ & $C$ & $10^{-7}$ & 18\\
$h-10^{-3}$ & $6\kappa_*$ & $10^{-5}$ & 6\\
$1/2$ & $6\kappa_*$ & $10^{-3}$ & 1\\
\bottomrule
\end{tabular}
\end{center}
The lower interval therefore uses the old estimate. At the junction,
$C<6\kappa_*$ and $2y-d>0$, so the local exponent is no larger
than the old one. Affinity bounds the intermediate local rows. Above
$h$, a scalar slope bound of $2$ leaves saving
$10^{-7}-2\cdot10^{-10}>9\cdot10^{-8}$.

For larger free shifts the source's existing negative coefficient
contribution must be kept. Write $q=\delta p$, so $0\le q\le\delta/2$.
The audit verifies, for the normalized available width in this window,
\[
 -\frac{\delta t}{2}-\frac{qz_{\rm old}(t)}2
 \le-\frac{qz_{\rm new}(0)}2,\qquad t\ge10^{-3}.
\]
Here $z_{\rm old}(t)=\max(0,b-2t)/c$ with
$6\kappa_*\le c\le9/2$, and $z_{\rm new}(0)=b/C$; $b$ is the
residual budget after the unshifted reflected plain factors. The
normalized reserve at the first shift exceeds $3/10000$ and increases
thereafter, including after the old available width becomes zero.
The free-shift comparisons are exact scalar interface checks.
Section~\ref{sec:physical-scale-shortening} explains why downward changes
of the reflected plain scales can use either the local moment or the old
moment without changing the weighted coefficient mass. If the selected
physical prime product is shorter than the local lower bound, the old
moment has reserve $1-2(211/500)-(9/2)(3/100)=21/1000$.
The slot-selection theorem still measures its gain at the nominal target
capacity. Section~\ref{sec:endpoint-assembly} proves this distinction and
covers every physical Euler term, including an empty selected product.

\subsection{The noncircular feedback mechanism}
Section~\ref{sec:endpoint-assembly} inserts Proposition~\ref{hyp:local}
into the original weighted numerator with its coefficients, masks,
conductor losses, principal terms, and family reductions preserved.
The numerical feedback used in that proof is as follows. Let $\beta_*$ be the supremum used in the source, with its sentinel
$1/2$. The $7/8$ initialization gives $\beta_*\le7/8$. If
$\beta_*>\theta$, put $\kappa=2\beta_*-1$; then
$\kappa_*<\kappa\le3/4$.

For the new upper-row expression, only the detector depends on $\kappa$.
The source bounds $\partial_\kappa c_\kappa<1/18$ and
$\partial_cR\le5/12$ give
\[
 \partial_\kappa F_{C,\kappa}
 \le\frac{3h}{4}\frac5{12}\frac1{18}
 =\frac{3969737987}{288000000000}<\frac1{50}.
\]
Consequently Proposition~\ref{prop:certificate} implies
\begin{equation}\label{eq:feedback}
 F_{C,2\beta_*-1}(\delta,p)-\beta_*
 \le-10^{-7}-\frac{24}{25}(\beta_*-\theta)<0.
\end{equation}
For old rows the coefficient $6\kappa$ also varies. The separate
derivative estimate in Section~\ref{sec:endpoint-assembly} gives
$\partial_\kappa G_o<1/8$, leaving at least
$3(\beta_*-\theta)/4$ in addition to their certificate savings.
Together these strict savings feed the source's contradiction to the
definition of $\beta_*$. The hypothesis on the slot zero bound is an
equality at the chosen $\kappa$; the improved zero-free conclusion is
not used as an input. The final supplement supplies the analytic
numerator and continuation steps that turn this feedback into the
proposed conclusion $\beta_*\le\theta$.

\section{What is proved in writing and what is independently checked}\label{sec:gaps}

Draft 3 supplies a complete proposed written analytic argument for the
local moment and the all-family endpoint, relative to the pinned source
results named in its supplements. It addresses the questions left open
in Draft 2: signed characters and zero extensions, the fixed dictionary
and unmarked plain moment, uniform height thresholds, canonical low-row
restriction, enlarged radii, and the full local branch ledger. The final
assembly includes shortened plain factors, a short or empty selected
prime product, free Euler shifts, principal and tail terms, continuation,
and the passage to ordinary Dirichlet and zeta functions. These arguments
are supplied for review; their inclusion does not create a verified
tracker record.

The established finite evidence has a narrower scope. The rational audits
cover continuous sign inequalities, exact symbolic exponents, and complete
rational Bernstein certificates. Finite CRT and joint-fiber checks help
detect mistakes but are not used in place of the general reconstruction
proof. The current development replay compiles 15 components and an audit
module, with 48 printed theorem axiom checks. They use cached pinned
dependencies, so this is not a fresh rebuild of the entire dependency graph.

\begin{center}\small
\begin{tabularx}{\textwidth}{@{}>{\raggedright\arraybackslash}p{0.31\textwidth}X@{}}
\toprule
Evidence group & What it establishes\\
\midrule
Finite deletion and sixth powers & Preservation of natural holes and positive-half-plane zero data.\\
Buffered slots and restricted incumbent & Pointwise payment on selected rows, then application of the old remaining-state moment.\\
Actual unmarked inputs and plain adapter & The fixed-family plain and inverse input estimates, with actual witness and coefficient dictionaries.\\
Height-uniform density & Uniform thresholds, bounded windows and the sixth-power lift, using its stated source inputs.\\
Signed and canonical-source lemmas & Opposite twist signs, complete mask maxima, fixed-modulus nonprincipality, actual restricted coefficient factorization and the padded-radius ledger.\\
Scale-shortening lemma & The arithmetic dichotomy paying sufficiently shortened plains by the old moment.\\
Endpoint certificate & A scalar saving below $\theta$ on the complete parameter rectangle.\\
\bottomrule
\end{tabularx}
\end{center}
The 48 reported axiom lists contain only \code{propext},
\code{Classical.choice}, and \code{Quot.sound}. This does not constitute
an axiom audit of a new complete endpoint theorem. The package manifest,
commands, input hashes and receipts identify the exact component scope.

Independent mathematical review of the complete composition is still
needed, including its use of the pinned incumbent analytic results and
its order of quantifiers. No new complete all-family Lean theorem has
been constructed. Clean dependency replay, comparison against the
canonical target statements, and independent kernel verification have
not been completed. These are validation limitations, separate from
the written proof now supplied. Until that review and verification are
complete, this manuscript should be described as a proposed analytic
proof under review, not as an independently certified improvement.

\section{Attribution and source relationship}\label{sec:credits}

OpenAI's development supplies the initial $7/8$ theorem and the analytic
and arithmetic framework. Levy's work supplies the weighted-numerator
input adapted by the Cycle25 source. Collet's work supplies the variable
moment, detector, quartic optimization, and source architecture used
throughout this draft. Our proposed contribution is the local
$4.498$ refinement, its joint-label analysis, and its proposed endpoint
assembly supported by an exact rational calculation. These descriptions do not imply endorsement by
any upstream author or maintainer.

This manuscript develops the three candidate notes in the research
package \cite{gao}, the second-revision input and reconstruction work
\cite{gaov2}, and the new source arguments included in the supplements.
It distinguishes their mathematical claims from their current validation
status. Research,
code, formalization, and writing received AI assistance; the internal
reviews recorded with the package are model-assisted development
reviews, not independent human referee reports. The package retains its
Apache-2.0 license and upstream notices.

\appendix
\section{Reproduction and the certificate format}\label{app:reproduce}

The package linked in \cite{gao} contains three notes under the directory
\path{research/framework/programs/fixed_strip/candidates/}:
\begin{itemize}
\item \code{LOCAL\_PHYSICAL\_MOMENT\_EXTENSION.md};
\item \code{JOINT\_HIGH\_ZERO\_SQUARED\_LABEL\_PAYMENT.md};
\item \code{LOCAL\_MOMENT\_TRACKER\_REFINEMENT.md}.
\end{itemize}
The first explains the proposed local moment, the second isolates a
useful squared-label branch, and the third records the endpoint
geometry. In particular, the coefficient $21/5$ used illustratively in
the second note is not the coefficient claimed in
Proposition~\ref{hyp:local}, and is not an established stronger moment.

From the first snapshot root
\path{proofs/gao-local-moment-20261011}, with Python and the recorded
SymPy version available, run the following commands. These paths belong
to that original snapshot, not to the present revision directory.
\begin{verbatim}
python3 reproduce.py --fetch-pinned-paper
python3 research/tools/local_physical_moment_audit.py --check
python3 research/tools/local_moment_endpoint_audit.py --check
\end{verbatim}
The fetched manuscript is a pinned input, checked by its stored digest;
it is not a paper authored by this project. The outputs in
\code{research/numerics/results/} record the exact assertions and their
scope. Both reports explicitly deny verification of the new full
analytic theorem. The present revision also supplies the following additional exact check:
\begin{verbatim}
python3 endpoint_extra_arithmetic.py --check
\end{verbatim}
It verifies 20 rational comparisons used by the written endpoint assembly.
The current component replay is documented with
\path{replay_components.py} in the revision README; it has separate
compiler and pinned dependency prerequisites.

For completeness, the Bernstein certificate used in
Proposition~\ref{prop:certificate} is based on the identity
\[
 f(u,v)=\sum_{i=0}^{m}\sum_{j=0}^{n}
 b_{ij}\binom mi u^i(1-u)^{m-i}\binom nj v^j(1-v)^{n-j},
 \qquad 0\le u,v\le1.
\]
The basis functions are nonnegative and their sum is one. Thus
$b_{ij}\ge0$ for all $i,j$ proves $f\ge0$ on the whole square.
An exact affine substitution maps each recorded rational rectangle to
that square. Rational reconstruction of the coefficients and exact
coverage checks make the conclusion inspectable without trusting
floating-point optimization.

\clearpage
\section*{Guide to the supplementary source arguments}
\addcontentsline{toc}{section}{Guide to the supplementary source arguments}
The following sections are part of the proposed analytic proof. They retain
the exact pinned source interfaces so a reader can inspect each dependency.
The local and endpoint assemblies are proposed proofs under review.
Their distinction from independent certification is stated in
Section~\ref{sec:gaps}.

\section{Signed rows, zero extensions, and label reconstruction}
\label{sec:signed-reconstruction}

The purpose of the reconstruction is to identify one integer row to which
the unmarked density theorem applies. It must identify the character on
all arguments, including its zeros, and must also explain which arithmetic
labels can be recovered from that row. Merely identifying a primitive
character on coprime arguments would lose the natural masks.

\subsection{The two signs at a common prime}
For supported primary elements define
$c_n(v)=\mathfrak s_v((n))$, the source's zero-extended sextic ideal row.
Let $q(\pi,n)$ be the fixed reciprocity phase, so
\[
 c_\pi(n)=q(\pi,n)c_n(\pi).
\]
The phase belongs to the fixed ray dictionary. If $\pi\mid n$, both
sides vanish; the identity is not restricted to the unit case.
For a common prime with active net exponent $1\le t\le5$, it follows that
\begin{align*}
 c_n(\pi)c_\pi(n)^t
   &=q(\pi,n)^t c_n(\pi)^{1+t},\\
 c_n(\pi)\overline{c_\pi(n)^t}
   &=q(\pi,n)^{6-t}c_n(\pi)^{7-t}.
\end{align*}
On units this is ordinary multiplication of sixth roots of unity;
off units each displayed product is zero. Thus the left coefficient has
residue $1+t$ and the right coefficient residue $1-t$, modulo six.
These are the signs in the actual formal coefficients
\path{leftCoefficient} and \path{rightCoefficient}.
A manuscript convention assigning the conjugate coefficient to the
C-side must exchange the side names consistently.

If $t=5$ on the left, or $t=1$ on the right, the resulting exponent is six.
The source's exact identity is
\[
 c_n(\pi)^6=\mathbf1_{(n,\pi)=1}.
\]
Consequently reducing the exponent to zero moves this factor into the
puncture; it does not replace it by one. Multiplying the primewise
identities proves the corresponding complete first-row character
identity with every canceled prime still punctured.

For one side let $F=\prod_{\pi\mid r}\pi^{e(\pi)}$, choosing the
source's primary prime generators, and let $L=\log_Z NF$. The ideal
used by the formal source has generator
$\operatorname{idealGenerator}(\mathfrak r)$, which can differ from the
product of chosen generators by a unit. Retain that unit in the effective
row, or absorb its character into the fixed six-unit dictionary. No
unproved equality of these two generator choices is needed.

\subsection{The complete second frequency}
For the second common-support ideals $D_2,E_2$, let $G$ be their common
gcd generator. The actual correlation vanishes unless $G\mid j$.
After writing $j=Gu$, the nonunit part of the equal-multiplicity
partition forces $V\mid u$, where $V$ is its radical generator.
Thus the genuine nonzero frequency is
\[
 j=GVh',\qquad h'\ne0.
\]
The source's partitioned correlation, divided by its stated normalizer,
has modulus at most one, including its extension by zero off the
partition. Multiplicativity now combines the first row with this frequency:
\[
 w=u_0\,eFGVh',
\]
where $u_0$ is the retained unit. Every surviving whole-product character
is a fixed ray factor times $c_n(w)$, with the common puncture and all
canceled-prime holes. Relative slot characters remain their original
fixed factors. Whole-character conjugation or the reflected row $-w$
is handled by whole-character height reflection and the fixed unit
list. The dictionary is a finite union of parent translates of the
allowed $Q_*$-character group. The buffered maximum is formed within
the actual parent translate; another translate is not added to that
maximum. Conjugating only its fixed factor would change this family
and is not an allowed replacement.

The exact source ingredients are
\path{SecondCanonicalFrequency.idealCorrelation_zero_outside_common},
\path{SecondCanonicalNonunit.nonunitFrequencyGenerator_dvd}, and
\path{SecondCanonicalFrequency.actual_ideal_partition_scalar}.
The new signed-row lemmas check the primewise identities and their
multiplicative combination against these source conventions.

With the actual common main ball of logarithmic radius
$K_0+(d-c)_++6a+\varpi$, the second frequency envelope gives
\[
 \log_Z Nw\le M-R-(c-d)_++L-6a-\varpi+\eta.
\]
Discarding the favorable $-\varpi$ gives \eqref{eq:integerrow}.
A nonempty integral row window always has $Nw\ge1$; negative nominal
logarithmic ranges are either empty or are filled by the reserved
bounded-window allowance. Density is not applied as a count smaller
than one on a nonempty window.

\subsection{Which labels are divisors of the row?}
The first residual projector label $s$ divides both second columns.
After the complete second common-support extraction it therefore divides
$G$, and hence divides $w$. The entire second common radical also divides
$G$. The source proves these assertions as
\path{SecondSourceMask.source_pair_divisor_common} and
\path{divisor_dvd_common}.

The second residual projector label $t$ is different. Its individual
terms belong to the artificial separated extension of the genuine
correlation, and $t$ need not divide $G$ or $w$. Freeze it while estimating
the row norm. Its signed priority allocation selects plain divisors and
freezes prime slots. If a selected plain is $n=dv$, then
\[
 \psi_w(dv)=\psi_w(d)\psi_w(v).
\]
The factor $\psi_w(d)$ is retained outside, even if it is zero. The
quotient $v$ may still be divisible by a prime of $d$; no new coprimality
condition is introduced. This explains why the primitive row and its
buffered bin survive allocation. Its divisor-dyad count is paid together
with the two exact normalization scalars in B3.67--B3.69.

Among the first common primes, those active with nonzero row residue
occur in $F$, and the selected mask-divisor primes occur in $e$.
The remaining first radical has logarithmic width at most
\[
 p-R-E+R_{\rm zero}.
\]
This accounts separately for inactive net-zero primes and for active
primes whose row phase cancels. They cannot be called divisors of $w$.
Once their radical is fixed, all remaining first labels are supported
on it together with $\operatorname{rad}(w)$.

For completeness, a uniform support-counting bound follows from Rankin's
inequality. If $NQ\le Z^B$ and $NI\le Z^T$, then for any $s>0$,
\[
 \#\{I:\operatorname{rad}I\mid Q,\ NI\le Z^T\}
 \le Z^{sT}\prod_{\mathfrak p\mid Q}(1-(N\mathfrak p)^{-s})^{-1}.
\]
Choose $sT$ small. Each sufficiently large prime factor in the product
is at most $(N\mathfrak p)^b$ for any prescribed $b>0$; the finitely
many remaining primes contribute a constant. Choosing $bB$ small proves
a bound $O(Z^\epsilon)$ uniformly in the support ideal. The field has
only $O(X)$ ideals of norm at most $X$, so counting the missing radical
costs its displayed norm exponent. Fixed finitely many slot assignments,
bounded-range powers and priority allocations have the same arbitrarily
small power allowance after the slot system is fixed.

This proves the general label reconstruction used in the high norm.
It retains the first labels until the joint row count, and it does not
charge the second projector count twice. Nonprincipality and primitive
conductor bounds are supplied by the following fixed-dictionary section;
they are not consequences of the size bound on $w$ alone.

% Written audit and bridge, 11 October 2026. Input fragment.
\section{The fixed dictionary behind the actual signed rows}
\label{sec:fixed-dictionary-bridge}

This section concerns the local family whose parent character and ray
modulus are fixed before the scale. It does not assert a uniform density
estimate for arbitrary parent characters with growing conductor. That
stronger assertion is unnecessary here: the moving factors in both
transformations are retained in the complete integer row.

Write $m_0=2(1-\omega)$. Let $Q$ be the \emph{literal} fixed ideal used
by the source's fixed-inducing-row predicate. Its hypotheses include
$Q\subseteq(72)$ and the stated containments in the original fixed ray
moduli. The relative slot characters have presentations of modulus
dividing $Q$. These containments use the ideal convention that
$Q\subseteq\mathfrak m$ means $\mathfrak m\mid Q$.
For the actual initial family $F$, this ideal is
\[
 Q=\operatorname{internalQ}
       (F.\mathrm{modulus}\cap(72))\,\eta_0
   =F.\mathrm{modulus}\cap(72)\cap\mathfrak q_{\eta_0},
\]
as defined by \path{NaturalFixedRaySourceFixedIdeal.sourceFixedIdeal}
and \path{sourceExceptional}. Thus the needed containments belong to
the literal source family; they are not a later enlargement of its
exceptional set. In a transformed child this ideal remains fixed even
though the separate whole-product character can move.

\paragraph{A finite collection chosen before arithmetic labels.}
Take the primary ray characters modulo $12$, the six unit-row characters
modulo $36$, and the allowed relative slot characters, together with their
inverses. On ideals prime to $Q$ their products belong to the finite group
of ray characters modulo $Q$. Multiply these products by each of the
original fixed parent characters; the resulting finite collection is a
union of fixed translates of that group. The parent itself need not have
modulus dividing $Q$. Different zero extensions at primes dividing $Q$
introduce only finitely many additional presentations. Thus there is a finite list
$\{\eta_\lambda:\lambda\in\Lambda\}$ containing every product that will
occur, chosen before the common supports, frequencies, or rows. Products
over an arbitrarily large number of active primes do not enlarge this
list: the reciprocity phases are characters of a single fixed residue
group. They do not introduce one new independent character for each prime.
Keep distinct translates of the $Q$-character group separate when defining
zero bins. The finite union is used to make constants simultaneous; it
does not imply that eligibility for one translate makes all other
translates nonprincipal. Partitioning the fixed outer labels by their
parent translate costs a fixed finite factor in the eventual joint count.

The source realizes the primary phases by
\path{FixedRay.primaryRayCharacter}, whose modulus is $(12)$, and unit
phases by \path{HeckeUnitRows.character}, whose modulus is $(36)$.
For the relative factors, the exact source formula is
\[
 \nu_\theta=\chi_\theta\eta_0^{-1},\qquad
 Q\subseteq M,\ Q\subseteq\mathfrak q_{\eta_0}
 \ \Longrightarrow\ Q\subseteq\mathfrak q_{\nu_\theta}.
\]
This is \path{NaturalFixedRaySourceCharacters.relativeCharacter_fixedQ}.
The allocation theorem \path{AllocatedRayDictionary.originalImage_fixedQ}
preserves it for every surviving slot. The full source identities
\path{both_common_rectangles} retain the original relative characters,
profiles and independent heights; allocation changes which original slots
remain, not their character dictionary.

\paragraph{Exact coefficient identity.}
Use the signs and primary-generator conventions of the signed-row
reconstruction. For either first-side orientation the complete second
integer row is
\[
 w=u\,eF_{\pm}GVh\ne0,
\]
where the fixed unit $u$ may instead be absorbed into the finite list.
Every actual residual character coefficient, including a surviving
relative slot twist, is of the form
\[
 c_{\eta_\lambda}(I)\,\mathfrak s_{m_0^6w}(I)
       \mathbf1_{(I,\mathfrak R)=1}.                 \tag{FD}
\]
Here $\mathfrak R$ is the retained natural and common puncture. This is an
identity on all ideals, including the zero values. On supported ideals
it follows by evaluating the signed identities at the primary generator;
on unsupported ideals both sides vanish because of $m_0$. A canceled
exponent at an active prime produces a factor
$\mathbf1_{(I,\mathfrak p)=1}$ and therefore belongs to $\mathfrak R$.
The second extraction contributes $GVh$ multiplicatively to the same
integer row. The complete second coprimality projector is expanded with
its original signs. Extracting one of its assigned divisors $d$ uses
$\psi(dn)=\psi(d)\psi(n)$; it does not impose an additional condition
$(d,n)=1$. A zero outside coefficient annihilates that term. These facts
are exactly the distinctions made in the signed-row reconstruction.

The two first-side orientations can be treated separately, at a fixed
factor of two. Within the square of one first-side column the reflected
second rectangle has the row $-w$, which is covered by the same finite
unit-character list. There is no claim that $F_+$ and $F_-$ are equal.
If a whole rectangle is conjugated, retain that conjugation of the whole
character: its zero maximum is unchanged after reflection of the height
coordinate. Do not replace $\overline{\eta\,\mathfrak s_w}$ by
$\overline\eta\,\mathfrak s_w$. Such a replacement could change the
parent translate and its exceptional rows.

Choose the natural character $\chi_{\lambda,w}$ for the unmasked part of
(FD). The natural constructor gives
\[
 N\mathfrak q_{\chi_{\lambda,w}}
 \le N\mathfrak q_{\eta_\lambda}\,N(m_0)N(72)\,Nw
 \le K Nw,                                             \tag{FC}
\]
with one fixed $K$. Its ideal coefficients equal those of
\texttt{momentData $\eta_\lambda$}; the new theorem
\path{natural_row_moment_dictionary} verifies this equality.
One must use this natural representative for (FC), since the abstract
choice inside \path{RowData.character} has no stated conductor bound.
The primitive inducing character has still smaller conductor.

\paragraph{Literal exclusion survives the change of presentation.}
For a nonzero mask and row, the source proves
\[
 \operatorname{FI}_Q(\eta;m_0r,1,w)
 \ \Longleftrightarrow\
 \operatorname{FI}_Q(\eta;m_0,1,w),                       \tag{FE1}
\]
by \path{FixedRowMask.fixedInducingRow_mul_mask_iff}.
If $Q\subseteq\mathfrak q_\nu$, it also proves
\[
 \operatorname{FI}_Q(\eta\nu;m_0,1,w)
 \ \Longleftrightarrow\
 \operatorname{FI}_Q(\eta;m_0,1,w),                       \tag{FE2}
\]
by \path{FixedRayInducingTransport.fixedInducingRow_product_iff}.
The reverse implication is not obtained by canceling a possibly zero
character value. Multiplication by $\nu^{-1}$ recovers the original
character with the additional modulus mask of $\nu$; (FE1) then removes
that mask from the \emph{inducing predicate}. This is why the equivalence
remains valid on the literal zero-extended family.

Apply (FE1), (FE2), and (FD) to the actual source exclusion. They identify
it with the same predicate on $\chi_{\lambda,w}$. In particular every
allowed relative twist in this same parent translate is nonprincipal on
an eligible row. No conclusion is made here about another parent
translate. The source
states this consequence directly as
\path{NaturalFixedRaySourceCharacters.sector_nonprincipal}.
For the second common extraction,
\path{SecondExceptionalFamily.Family.inducing_iff} and
\path{paired_inducing_iff} implement precisely the reduced-row and
reflected-row versions of this transport. The ideal $Q$ is not enlarged.
It would be incorrect simply to replace the source exclusion by
nonprincipality of one preferred primitive representative.

\paragraph{Zero bins and the density input.}
For one fixed parent translate, take the maximum of the zero real parts
over its allowed $Q$-character twists, and use the buffered height boxes
of the source. Whole-character inverses have the same maximum and need
not be modeled by an additional translate. Each actual character in
that maximum is a masked presentation
of one of (FD). Finite Euler factors are nonzero in $\Re s>0$, so its
zeros in these boxes equal those of its primitive inducing character.
Removing sixth powers, $w=ub^6$, introduces further Euler holes and
does not change these positive-real zeros. The fixed-dictionary density
proposition therefore applies to the selected free core $u$, followed
by the height-uniform sixth-power lift.

If a bin is selected by the maximum over all allowed labels in this translate,
some label attains that maximum. The union over this \emph{fixed} finite
list costs only a constant. The upper side of the buffered bin bounds
every surviving prime-slot character, while the lower side supplies the
zero used for density. This step neither sums over the moving arithmetic
labels nor freezes them before counting. Those labels are recovered from
$w$ and the missing radical by the signed-row reconstruction.

Consequently, for $7/10\le\delta\le3/4$, the actual eligible integer rows
in the selected high bin satisfy
\[
 \#\{w:Nw\le H\}\ll_{\varepsilon}
       (1+T)^J H^{\rho(\delta)+\varepsilon},\qquad
 \rho(\delta)=1-\frac{10\delta}{3(5-2\delta)},
\]
up to the stated arbitrarily small bin-buffer loss. The constants and
threshold precede $T$, every live common-support label, and $H$.
Polynomial-size masks do not enter the dictionary size or the primitive
conductor in (FC); their smooth estimates are paid separately by uniform
Euler deletion. The pointwise upper bounds use the common buffered bin,
whereas the count only needs the label attaining its lower edge.

\paragraph{Fourier kernels whose parameter varies with labels.}
The common Fourier argument also permits a scale parameter $A_\ell$
depending on a retained label $\ell$. Uniform Euler derivative bounds
give Fourier densities satisfying
$|f_A(t)|\le C(1+|t|)^{-2r}$, uniformly in $A$. Choose once and for all
$d\mu(t)=C(1+|t|)^{-2r}\,dt$ and write
$f_{A_\ell}(t)\,dt=a_\ell(t)d\mu(t)$, where $|a_\ell(t)|\le1$.
The integral of the joint vector then contains the diagonal contraction
$\operatorname{diag}(a_\ell(t))$. Its operator norm is at most one,
so Minkowski's inequality gives exactly the same joint bound, with no
extra label count. This avoids the unjustified assertion that the
individual Fourier measures for distinct $A_\ell$ are equal.

This establishes the fixed-dictionary, conductor, mask, and literal
eligibility bridge from the actual signed-row identities. The remaining
assembly of the physical moment must still verify its normalization,
comparison and exceptional terms, and its use of the reserved exponent
margins. Those conclusions do not follow from this bridge alone.

% Research derivation. This file is an input fragment, not a standalone paper.
% The plain-moment theorem below has a compiled Lean counterpart.
% The final assembly of the zero-density statement is a written proof.
\section{The unmarked plain moment and density for a fixed character dictionary}
\label{sec:actual-unmarked-density}

This section closes the zero-slot plain-moment input for a fixed finite
character dictionary. It does not identify the varying rows in the proposed
physical moment with such a dictionary; that separate identification must
preserve every Euler hole and the stated constants.

Let $\mathcal O=\mathbb Z[\omega]$ be the Eisenstein integers and put
$m_0=2(1-\omega)$. A free row means a nonzero element $u$ whose principal ideal
has every prime valuation less than six. Fix a nonempty finite set $\Lambda$ (the empty case is immediate) and
characters $\eta_\lambda$, $\lambda\in\Lambda$. For each $u$ choose the
\emph{natural} row character $\chi_{u,\lambda}$, so that
\[
 c_{\chi_{u,\lambda}}(I)
 =c_{\eta_\lambda}(I)\,\mathfrak s_{m_0^6u}(I),
 \qquad
 \mathrm N\mathfrak q_{u,\lambda}
 \le K_\lambda\,\mathrm N(u),
 \quad K_\lambda=\mathrm N\mathfrak q_{\eta_\lambda}
                   \mathrm N(m_0)\mathrm N(72).
\]
Here $\mathfrak s$ is the literal extended sextic ideal character, including
its zero values at ramified primes. The conductor estimate refers to the
natural representative obtained from the conductor-bounded row constructor.
The abstract \path{RowData.character} is chosen by a different existence
theorem and does not carry this bound. Its coefficients agree with the
natural representative, which is all the moment estimates require.
Whole-character conjugation preserves the modulus and reflects the zero
height. It can be handled through this symmetry; it must not be confused
with conjugating only the fixed factor $\eta_\lambda$.

\subsection{The zero-slot plain input}

For $n\le2$, $\sigma\in[0,1]$, $t\in\mathbb R$, write
\[
 P_n(\chi;D,\sigma,t)
 =D^{-1/2}\sum_{I\ne0}c_\chi(I)
   (\log(\mathrm NI/D))^n W(\mathrm NI/D)
   (\mathrm NI/D)^{-\sigma+it},
\]
where $W$ is the fixed positive annular detector profile, supported in
$[1/4,9/4]$. Its logarithmic derivatives occurring here are those of the
existing detector construction.

\paragraph{Plain-moment proposition.}
Assume the contradiction branch $\beta\ge51/100$. For each $\epsilon>0$
there are $C>0$, $J\in\mathbb N$, and $U_0$, depending on the fixed
dictionary and $\epsilon$, such that for every $U\ge U_0$, every height
$H\ge0$, every unmarked detector fiber whose right length satisfies $m\le3/4$,
and whose finite row set $\mathcal R$ satisfies $\mathrm N(u)\le U$ and
the displayed coefficient dictionary, one has
\[
 \sum_{u\in\mathcal R}
 |P_j(\chi_{u,\lambda};U^m,\sigma,t)
  P_k(\chi_{u,\lambda};U^m,\sigma,t)|^2
 \le C(1+H)^J U^{\max(1,2m)+\epsilon},
 \tag{P}
\]
for $j+k\le2$, $\sigma\in[0,1]$ and $|t|\le H$.
The threshold precedes the height quantifier. No lower bound on
$\mathrm N(u)$ and no prime-slot supply is assumed. In the formal theorem
$\mathcal R$ is the row set of an unmarked detector fiber; the proof of its
plain estimate only uses its row bound, coefficient dictionary, and $m$.

\paragraph{Proof.}
Fix a label and choose, for example,
$Q=\mathfrak q_\eta(72)$. Then $Q\ne0$, $Q\ne\mathcal O$,
$Q\subseteq(72)$, and $Q\subseteq\mathfrak q_\eta$. These are the exact
fixed-modulus gates of the incumbent analytic induction. For
$d=\min(\epsilon/4,1/4)>0$, its terminal certificate with
$\kappa=\max(3/4,2\beta-1)$ supplies the zero-slot bound at error $d$.
The zero slot profile is a legitimate input to that certificate; no
improvement of the physical moment is used. The fixed character modulus
satisfies $\mathrm N\mathfrak q_\eta\le U^d$ for all sufficiently large $U$.
The natural initial state has width
\[
 1+d+\max(0,2m-1)=\max(1,2m)+d\le2.
\]
Consequently the incumbent zero-slot estimate, followed by the uniform
Schwartz-profile bounds, proves (P) on rows satisfying
\path{initialKeep}, with exponent $\max(1,2m)+2d$.
The profile control contributes $C(1+H)^J$ and does not change the scale
threshold. The legacy proof interface is reached by adding one auxiliary
slot of width $7/37$ with identically zero profile and always selecting the
empty product. Every legacy structure field, including its numerical
supply field, is then proved; no prime sum or supply property is imposed
on the original family.

The omitted free rows form a fixed finite set $E_\eta$. Indeed, the
fixed-inducing-row theorem puts all their prime factors in the finite
support of $Q$, and free-row valuations are at most five. There are only
finitely many resulting ideals and six possible units per principal ideal.
This is a fixed finite set, not the whole interval of small rows.

The exceptional contribution has the required size without cancellation.
The detector Schwartz profiles have a uniform absolute bound $A$ for all
$n\le2$, $\sigma\in[0,1]$, and all frequencies. Since
$\#\{I:0<\mathrm NI\le x\}\le128x$, coefficient norms are at most one,
and the profile vanishes above $9/4$, the normalized finite sum gives
\[
 |P_n(\chi;D,\sigma,t)|\le288A\sqrt D
\]
for every character, including principal characters. Therefore
\[
 \sum_{u\in\mathcal R\cap E_\eta}|P_jP_k|^2
 \le |E_\eta|(288A)^4D^2
 \le |E_\eta|(288A)^4U^{\max(1,2m)+\epsilon}.
\]
This proves (P) on all rows. Taking maxima of finitely many thresholds and
degrees, and a sum of the finitely many constants, proves the simultaneous
statement for $\Lambda$. The compiled theorem
\path{fixed_unmarked_plain} proves this argument for one label; finite
label closure is an ordinary finite maximum argument.

\subsection{From actual witnesses to a height-uniform density bound}

Let $N(Y,a,T)$ count free rows $u$, $\mathrm N(u)\le Y$, for which at least
one dictionary character has a zero with real part at least $a$ and
imaginary part of absolute value at most $T$. In this paragraph
$17/20\le a\le7/8$. Use the incumbent bound $\beta\le7/8$. If $\beta<a$, the
nonexceptional contribution is zero; the fixed exceptional set is handled
by its constant cardinality. The lower hypothesis $\beta\ge51/100$
in (P) is automatic whenever the nonexceptional contribution is nonempty
and hence is not an extra theorem assumption on the desired zero-free result.

\paragraph{Density proposition.}
For every $\varepsilon>0$ there are $C,J,Y_0$, fixed before $a$ and $T$,
such that for all $Y\ge Y_0$ and $T\ge0$,
\[
 N(Y,a,T)\le C(1+T)^JY^{\rho(2a-1)+\varepsilon},
 \qquad
 \rho(\delta)=1-\frac{10\delta}{3(5-2\delta)}.
 \tag{D}
\]
The constant depends on the fixed dictionary. The same assertion for a
scale-dependent dictionary does not follow from this proposition.

\paragraph{Proof.}
Set $K=\max(1,\max_\lambda K_\lambda)$ and $U=KY$. Then both the row norm
and every natural character modulus are at most $U$. This is a fixed
multiplicative padding; no power of $Y$ is lost. Remove the finite union
$E=\bigcup_\lambda E_{\eta_\lambda}$. On the remaining rows every natural
character is nonprincipal, by the incumbent nonprincipal gate. The finite
union contributes $O(1)$ to (D), so there is no need to construct detector
witnesses for its members.

Choose small detector and moment errors before the scale parameter. The
actual inverse moment for each fixed \path{momentData} family and (P)
now give both moment inputs with one bound $C(1+H)^{J_0}$. Both the constant
and degree are independent of the detector zero bin, the dyadic indices,
and $U$. The actual detector lengths satisfy
$m\le1/2+75\epsilon_{\rm det}\le23/40$, so (P) applies.
Choose $\tau_{\rm det}>0$ \emph{after} $J_0$, sufficiently small for the supported
witness construction and for
$\tau_{\rm det}(J_0+1)$ to fit within the allotted final error.

First suppose $T\le U^{\tau_{\rm det}}$. For each row, apply the existing finite
buffered-bin lemma to its finite character family with height step
$U^{\tau_{\rm det}}>2$ and mesh $e>0$. It gives an integer $1\le i\le I$,
$I=\lceil49/(100e)\rceil$, and a real grid value $b=51/100+ek$ such that
\[
 b\le M(3iU^{\tau_{\rm det}})<b+e,
 \qquad M(3(i+1)U^{\tau_{\rm det}})<b+2e.
\]
Here $M$ is the actual maximum of zero real parts in the specified box,
with the usual $51/100$ fallback. Because the original row has a zero
with real part at least $a$ at height at most $T$, monotonicity gives
$b>a-e$. The incumbent upper bound gives $b\le7/8$. There are only
$O_e(1)$ possibilities for $(i,k)$, so partition the rows by these values.
For each part set $\delta_b=2b-1$ and fix its balanced cutoff
$t_*=5/(5-2\delta_b)$ before constructing witnesses.

The actual supported-witness theorem, with base scale $Z=U$ and $d=1$,
produces witnesses in each nonempty part. Its hypotheses are now literal:
nonprincipality, modulus at most $U$, the two buffered maxima, and the
natural-row coefficient identity. It uses frequency allowance
$U^{\tau_{\rm det}/2}$, and the witness frequency is bounded by
\[
 H_i=2\pi U^{\tau_{\rm det}/2}+3iU^{\tau_{\rm det}}\ll_e U^{\tau_{\rm det}}.
\]
Choose the auxiliary small parameters in the stated detector budget
$12e(22+2)+8\kappa_{\rm aux}+2\eta_{\rm aux}
 \le\epsilon_{\rm det}/2$ and
$4\tau_{\rm det}<\eta_{\rm aux}$. They can all be chosen positively, with
$e<1/1000$, and then $\tau_{\rm det}$ decreased further as necessary. No height
variable occurs in the resulting eventual threshold.

Partition once more by the witness label and its two dyadic indices, all
constructed using this fixed $t_*$. The no-slot inverse and plain estimates yield the
already compiled balanced count
\[
 |\mathcal R_{\rm fiber}|
 \ll (1+H_i)^{J_0+1}
 U^{\rho(\delta_b)+156\epsilon_{\rm det}+\epsilon_m}.
\]
The fixed label count and the two dyadic counts cost an arbitrary positive
power of $U$. Since $b>a-e$ and $\rho$ is decreasing, differentiation on
the relevant compact interval gives
$\rho(2b-1)\le\rho(2a-1)+3e$.
Choose $\epsilon_{\rm det},\epsilon_m,e$ and the dyadic padding so their
sum with $\tau_{\rm det}(J_0+1)$ is less than $\varepsilon$. Summing the finite
partitions proves $N(Y,a,T)\ll U^{\rho(2a-1)+\varepsilon}$ in this case.
The replacement $U=KY$ only changes the fixed constant.

If $T>U^{\tau_{\rm det}}$, the elementary lattice count is enough:
$N(Y,a,T)\ll U$. Since $\rho(2a-1)>0$ and $U\ge1$, for any integer
$J\ge1/\tau_{\rm det}$,
\[
 U\le T^J\le(1+T)^JU^{\rho(2a-1)+\varepsilon}.
\]
This proves (D) for arbitrary $T$, with a common threshold fixed before
that height. Bounded $Y$ are covered by enlarging the constant using the
finite row ball, uniformly in $T$ and $a$.

When this result is used at an external physical height $T=Z^{\tau_{\rm out}}$,
choose $\tau_{\rm out}$ only after the final degree $J$ in (D) has been
fixed, so that $J\tau_{\rm out}$ fits the physical error allowance. The
external exponent $\tau_{\rm out}$ is distinct from the internal detector
exponent $\tau_{\rm det}$ chosen in this proof. Reusing a single exponent
for both purposes would make the parameter selection circular.

The exponent satisfies $\rho(\delta)\ge2/7>1/6$ for
$\delta\le3/4$. Thus the weighted sixth-power raw-row lifting proved in
\path{HeightUniformDensity.lean} applies to (D). This completes the
analytic unmarked density input for a fixed dictionary. Applying it to the
joint physical row still requires the exact signed coefficient identity,
a fixed-dictionary reduction for every varying label, and the stated
conductor and mask accounting.

\subsection{Formal verification scope}

The new module \path{PlainUnmarkedAdapter.lean} contains nine audited
objects: the natural-row coefficient dictionary, the legacy embedding, the retained estimate, the actual zero-slot
input, the retained eventual field, finiteness of exceptional free rows,
the absolute detector-polynomial bound, the finite-row absolute bound, and
the complete eventual plain field. All nine printed axiom sets contain
only \path{propext}, \path{Classical.choice}, and \path{Quot.sound}.
Compilation reuses development dependency objects; this is not a clean
source replay or an independent checker certificate. The assembly of (D)
above is a written mathematical proof, not an additional Lean theorem in
this module.

% Development supplement, 11 October 2026.
% These are analytic transport lemmas, not a proof of the new local moment.
\section{Uniform deletion and common Fourier transport}
\label{sec:uniform-transport}

The following estimates resolve two bookkeeping questions in the proposed
argument. They neither identify a transformed character nor supply a
missing mean-value estimate. The field throughout is the fixed Eisenstein
field. Constants may depend on the fixed field, the indicated compact
supports, and the requested error exponent.

\begin{proposition}[Uniform finite Euler deletion]
Fix $B\geq 0$, $\sigma_0>0$, an integer $j\geq0$, and $\epsilon>0$.
Let $Z\geq2$, let $\mathfrak q$ be a squarefree integral ideal with
$N\mathfrak q\leq Z^B$, and choose complex numbers
$|a_{\mathfrak p}|\leq1$. Put
\[
 E_{\mathfrak q}(s)=\prod_{\mathfrak p\mid\mathfrak q}
       (1-a_{\mathfrak p}(N\mathfrak p)^{-s}).
\]
Then $E_{\mathfrak q}$ is nonzero on $\Re s>0$, and, uniformly in
$\mathfrak q$, the coefficients, and $\Im s$,
\[
 |E_{\mathfrak q}^{(j)}(s)|+
 |(E_{\mathfrak q}^{-1})^{(j)}(s)|
 \ll_{B,\sigma_0,j,\epsilon} Z^\epsilon
 \qquad(\Re s\geq\sigma_0).
\]
\end{proposition}
\begin{proof}
Every factor is nonzero because
$|a_{\mathfrak p}(N\mathfrak p)^{-s}|<1$ in the positive half-plane.
For $\Re s\geq\sigma_0/2$, both the product and its reciprocal are
bounded by
\[
 \prod_{\mathfrak p\mid\mathfrak q}
       (1-(N\mathfrak p)^{-\sigma_0/2})^{-1}.
\]
Given $b>0$, the factor at every sufficiently large prime norm is at
most $(N\mathfrak p)^b$: its logarithm is
$O((N\mathfrak p)^{-\sigma_0/2})=o(\log N\mathfrak p)$.
The finitely many smaller primes contribute a constant, independently
of the chosen ideal and coefficients. The product is consequently
$O((N\mathfrak q)^b)$. If $B>0$, choose $bB\leq\epsilon$; if $B=0$
the product is empty. Cauchy's derivative estimate on disks of radius
$\sigma_0/2$ proves the claim for every fixed $j$.
\end{proof}

Thus, once the exact identity
$L_{\rm masked}(s)=L_{\rm primitive}(s)E_{\mathfrak q}(s)$ is
established, extra polynomial-size holes preserve its positive-real
zeros and cost an arbitrarily small power in a contour argument staying
in a fixed positive half-plane. For a nonprincipal character the
identity also applies at $s=1$. A principal pole must still be treated
separately. This argument places no small-conductor hypothesis on the
\emph{masked presentation}. It applies the analytic estimate to the
primitive presentation and restores the finite factors afterward.
For the raw row family one still has to prove the primitive-conductor
bound and the exact coefficient identity; these are not consequences
of a product estimate.

For a prime polynomial with relative support $[c,C]\subset(0,\infty)$,
bounded coefficient modulus, and central normalization $P^{-1/2}$,
deleting primes dividing $\mathfrak q$ changes it by
\[
 O(P^{-1/2}\log(2N\mathfrak q))\qquad(P\geq1).
\]
Indeed there are at most $\log(N\mathfrak q)/\log2$ such primes.
Fixed powers of $\log N\mathfrak p$ add only powers of $\log Z$ on
polynomial-size scales. Real powers on the normalized support have
bounded modulus, and imaginary powers have modulus one. This error is
$O_\epsilon(Z^\epsilon)$ and is compatible with a squared pointwise
bound $P^\delta Z^\epsilon$ whenever $\delta\geq0$, after dividing
the error allowance between the original estimate and this correction.
No deletion of an arithmetic mask by positivity is involved.

\begin{proposition}[Fourier transport with labels inside the norm]
Fix a compact box $K\subset(0,\infty)^d$ and $J\geq0$. Suppose that
a family of smooth profiles $H_A$ is supported on one larger compact
box and has uniformly bounded Euler derivatives through order $2r$,
where $2r>J+d$. Then it has representations
\[
 H_A(y)=\int_{\mathbb R^d}\prod_{k=1}^d y_k^{it_k}\,d\mu_A(t),
 \qquad
 \int_{\mathbb R^d}(1+|t|)^J\,d|\mu_A|(t)\leq C_J,
\]
uniformly in $A$. The measure is chosen before any arithmetic label
or row is evaluated.

Let $\Lambda$ be a finite set of retained arithmetic labels, and let
$F_t=(F_t(\lambda,h))$ be a measurable family in
$\mathcal H=\bigoplus_{\lambda\in\Lambda}\ell^2(R_\lambda)$.
If
\[
 \|F_t\|_{\mathcal H}\leq B(1+|t|)^J,
\]
then
\[
 \left\|\int F_t\,d\mu_A(t)\right\|_{\mathcal H}
       \leq C_J B.
\]
In particular there is no factor $|\Lambda|$ in this transport.
The same conclusion holds for weighted direct sums with fixed
nonnegative arithmetic weights.
\end{proposition}
\begin{proof}
Pass to logarithmic coordinates $x_k=\log y_k$ and apply Fourier
inversion. Integrating by parts with $(1-\Delta)^r$ bounds the
Fourier density by $C(1+|t|)^{-2r}$. Compact support bounds the
necessary $L^1$ derivative norms uniformly. The asserted weighted
integrability follows because $2r>J+d$.

The second assertion is the triangle inequality for the Bochner
integral in the Hilbert space $\mathcal H$:
\[
 \left\|\int F_t\,d\mu_A(t)\right\|_{\mathcal H}
 \leq\int\|F_t\|_{\mathcal H}\,d|\mu_A|(t)
 \leq B\int(1+|t|)^J\,d|\mu_A|(t).
\]
The assumption supplies integrability. Alternatively, all row sums
may first be truncated, with the same estimate, followed by a limit.
\end{proof}

The required uniform derivative assertion holds for a kernel
$K(A\prod_k y_k^{e_k})$ if its Euler derivatives are bounded on
$(0,\infty)$ and the exponents $e_k$ are fixed. Explicitly,
\[
 D^\beta K\left(A\prod_k y_k^{e_k}\right)
   =\left(\prod_k e_k^{\beta_k}\right)
     (x\partial_x)^{|\beta|}K(x)
       \bigg|_{x=A\prod_k y_k^{e_k}}.
\]
Cutoffs on the fixed box and normalized norm powers introduce only
fixed seminorms and polynomials in their independent heights.
Relations among arithmetic norms restrict evaluation points inside
the box; they do not redefine $\mu_A$ for each label. Discontinuous
coprimality, valuation, or zero-bin indicators remain in the arithmetic
vector $F_t$ and are never differentiated.

The substantive premise is a bound on the \emph{joint} vector $F_t$.
Separate per-label bounds cannot be substituted for it without their
label count. Conversely, once an exact reconstruction bounds the joint
sum over $(\lambda,h)$, the Fourier step does not charge that label
count again. The preceding statement justifies this last passage but
does not prove the reconstruction or the required pointwise bounds.

There is also a useful version when an outer arithmetic label changes
the parameter $A$. The integration-by-parts proof gives a single
pointwise majorant for the Fourier densities,
$|f_A(t)|\leq C(1+|t|)^{-2r}$. Choose the positive measure
$d\mu(t)=C(1+|t|)^{-2r}\,dt$ before the labels. Write
$f_{A(\lambda)}(t)\,dt=a_\lambda(t)\,d\mu(t)$, where
$|a_\lambda(t)|\leq1$. For each $t$, multiplication by
$a_\lambda(t)$ is a diagonal contraction on the joint Hilbert space.
The same bound therefore follows with this one common dominating
measure; equality of the label-dependent Fourier densities is not
needed. This is used only for labels fixed outside the sum being
Poisson-transformed. It does not permit a discontinuous row-dependent
factor to be inserted into that sum before Poisson summation.

\paragraph{Quantifier order for height losses.}
If a density input is eventually bounded by
$C(1+T)^J H^\rho$, its threshold must be independent of $T$ before
one substitutes $T=Z^\tau$. An assertion proved only for each fixed
$T$ with an unspecified threshold $H_0(T)$ is insufficient. With a
uniform threshold, bounded $H$ are filled using the lattice count and
$(1+T)^J\geq1$, without a new height dependence. Fix the zero buffer,
the finitely many profiles, and their analytic height exponents first;
then choose $\tau$ so that the total $\tau J$ lies below the reserved
power saving. Frequencies outside the retained range require the
common Fourier tail estimate. The pointwise and density results must
be checked to have precisely this uniform order; it is not supplied
by the Fourier identity alone.

\section{Restricted canonical norms and the exact-width correction}
\label{sec:restricted-canonical}

This section explains why the low-row argument applies to the actual
allocated character sums. The key point is that the low indicator is a
restriction of a positive row norm. It never becomes a smooth profile,
a coefficient to be differentiated, or an input to Poisson summation.

\subsection{The actual coefficient and radial factorization}
Freeze a separated Fourier frequency and every arithmetic label outside
the current row sum, including the second projector label. After its
complete signed allocation, one positive rectangle norm is
\[
 \sum_z\mathbf1_{\mathrm{keep}(z)}\mathbf1_{\mathrm{low}(w(z))}
       \rho(z)\left|P_1(z)P_2(z)\prod_{i\in J}Q_i(z)\right|^2,
 \qquad \rho(z)\ge0.
\]
The fixed dictionary includes every relative sector in the current
parent translate. It gives a common bin for the slots in this expression.
The character is the exact one from Section~\ref{sec:signed-reconstruction}.

Let $r$ be the source's canonical radial state. Its keep predicate is
$z\ne0$ together with the literal exclusion for the moving character
$\tau$ and the fixed ideal $Q_*$. Define $r_{\rm low}$ by replacing only
that predicate with
$\mathrm{keep}(z)\wedge\mathrm{low}(w(z))$.
Its smooth radial profile, scale, whole-product character and punctures
remain the same. Its nonzero-row and nonexceptional conditions follow
immediately from those of $r$, so it is an admissible restricted natural
state of the same width.

The source identities apply to precisely this object:
\begin{enumerate}
\item \path{AllocatedNaturalRadial.common_allocated_normalized}
expresses the allocated row as the product of two natural plain sums and
the remaining natural prime slots. The two scale phases have modulus one;
common holes remain as the excluded-prime character.
\item \path{allocated_radial_bound} and its shared mask estimate
allow an arbitrary row keep predicate. Thus mask expansion can use
$r_{\rm low}$ and preserves the low restriction in every daughter.
Its plain scales are divided by their assigned ideals, and forced prime
slots are frozen.
\item \path{AllocatedChildren.commonChild_energy} identifies the
source's actual natural product with its allocated energy for any such
radial state. The new \path{actual_common_child_restricted} lemma
checks this identity on the canonical low rows.
\item \path{AllocatedRayDictionary.common_clipped_energy} reindexes
surviving slots to their original indices. It retains their relative
characters and replaces the frequency by the common shift plus the
original individual height. Its keep argument is arbitrary and therefore
still contains the low predicate.
\end{enumerate}
These are coefficient equalities before the moment estimate; arbitrary
absolute-value coefficients have not replaced the canonical slots.

For any nonnegative row weight, Cauchy applies with the same restriction
on both factors. Alternatively, one can first form a positive norm and
then split it into low and high parts. Neither operation asks for a
regularity property of the indicator. The smooth Fourier measure remains
outside the restricted arithmetic sums, as in
Section~\ref{sec:uniform-transport}.

\subsection{Payment before removing the restriction}
Choose one greedy subset $R\subseteq J$ from the actual clipped lengths.
On the low rows the buffered estimate gives, for each removed slot,
\[
 |Q_i(z)|^2\le B_i Z^{k_0z_i},\qquad k_0=7/10.
\]
The factors $B_i$ contain the allocated small losses and height allowance.
Multiplying these pointwise inequalities gives
\begin{align*}
 &\sum_z\mathbf1_{\mathrm{keep}(z)}\mathbf1_{\mathrm{low}(w(z))}
   \rho(z)|P_1(z)P_2(z)\prod_{i\in J}Q_i(z)|^2\\
 &\quad\le\Bigl(\prod_{i\in R}B_i\Bigr)Z^{k_0\sum_{i\in R}z_i}
   \sum_z\mathbf1_{\mathrm{keep}(z)}\mathbf1_{\mathrm{low}(w(z))}
   \rho(z)|P_1(z)P_2(z)\prod_{i\in J\setminus R}Q_i(z)|^2.
\end{align*}
Drop the low indicator only in the last nonnegative sum. It is now the
full canonical remaining-state norm of the old moment, with the same
moving character and fixed ideal. This is the exact factorization proved
by \path{restricted_energy_remove} and used by
\path{restricted_paid_incumbent}.

The two rates serve different roles. Removing one unit of slot width
reduces the old capacity defect by $9/2$, but its low-bin squared cost is
$7/10$. Greedy removal therefore costs at most
$(7/45)F_{\rm act}+(7/10)\eta_{\rm mesh}$.
If the slots run out first, the existing \path{ZeroAt} estimate applies
without a plain-capacity condition. There is no extra payment for a
residual plain defect. Bounded prime-slot lengths are handled directly;
the buffered theorem's large-length hypothesis is not applied at length
one. For a fixed finite slot list, per-slot losses are chosen from a
single prescribed aggregate budget.

\subsection{Actual width versus nominal width}
The canonical state uses the exact width
\[
 M_{\rm act}=\log_Z(\text{dyadic row scale})+
             \log_Z N\mathfrak q_\tau.
\]
It can be smaller than the nominal ledger width $M_{\rm nom}$. Put
$d=M_{\rm act}-M_{\rm nom}\le\eta_2$. The source's deleted-child capacity
estimate gives $F_{\rm act}\le\pos{F_{\rm nom}-d}$, so it is the combined
quantity $d+rF_{\rm act}$ which must be bounded. The elementary
inequality~\eqref{eq:widthpayment}, checked in
\path{width_plus_low_defect}, bounds it by
$\eta_2+rF_{\rm nom}$. Use $r=7/45$ on low main children and $r=1/6$
on the extracted errors.

For the padded actual main radius \eqref{eq:envelope}, direct
substitution gives
\[
 M_{\rm nom}=M-(c-d)_+-6a-\varpi-g_2+t_2
\]
and
\[
 (A-c-c_2-M_{\rm nom})-(A-M)
 =6a+\varpi-\min(c,d)-(c_2-g_2)-t_2\le6a+\varpi.
\]
These identities also hold with a fixed parent modulus width $q$, writing
$M=m+q$ and putting $m$ in the first nominal frequency formula.
They are checked in \path{RestrictedCanonicalSource.lean}.
The nominal low defect is consequently at most
$g+6a+\varpi+2\theta_N$ plus the fixed presentation padding. The source's
inner allowance is still $M_{\rm nom}+\Delta$, with
$\Delta=b_2-p_2+B_c+a\ge a$. The scale cancels in that subtraction.
Thus the strict reserve \eqref{eq:lowreserve} survives the exact-width
correction, with the padding chosen from it.

Finally, the proposed radius is already the formal common radius
\[
 \texttt{mainCommonRadius}
 =Z^{K_0+(d-c)_++2\sigma+\delta+\mathrm{reserve}}.
\]
Take $\sigma=3a$ and $\varpi=\delta+\mathrm{reserve}$, with
$2\le Z^{\mathrm{reserve}}$. The source's \path{main_energy_common}
proves smooth domination including the fixed support factor. The second
transform accepts every positive radial scale. The extracted errors keep
their original error radii and all three moving-character alternatives;
no low-bin heredity is used in those branches.

% Source-level analytic assembly, 11 October 2026.
% The proof below uses the pinned incumbent theorems and the new companion
% source derivations. It is not an independent verification of those imports.
\section{Assembly of the local physical moment}
\label{sec:physical-local-assembly}

The purpose of this section is to compose the source identities, rather than
to regard the favorable exponent ledger as an analytic estimate by itself.
All references B3.\(\cdot\) are to the pinned Cycle25 manuscript. We use its
existing $9/2$ moment and its unrestricted zero-slot moment, including
their literal natural coefficients and eligibility conditions. Neither is
replaced by the proposed stronger moment at any descendant state.
The unmarked density is the fixed-dictionary theorem proved in
\path{PlainUnmarkedDensity.tex}; its signed-row identity is supplied by
\path{SIGNED_ROW_RECONSTRUCTION.md}, and the actual fixed dictionary,
natural conductor, and literal eligibility transport are supplied by
\path{FixedDictionaryBridge.tex}. Uniform finite Euler deletion
and Fourier integration are supplied by \path{FourierMaskTransport.tex}.

\begin{proposition}[Fixed-family local physical moment]
Assume the incumbent bound $\beta_*\le7/8$, and work in its nontrivial
contradiction branch $\beta_*\ge51/100$. Fix an admissible natural
parent family of B3.4. Its whole-product character is a fixed finite
ray character times $\chi_n(k)$, with no parent moving support of
positive width. Fix its ray modulus, anchor character, and nonzero
fixed ideal. In particular retain the actual
$Q_* = \operatorname{internalQ}(Q,\eta_0)$ and the literal exclusion
$\neg\operatorname{FixedInducingRow}$ associated to that ideal;
$Q_*$ is proper, nonzero, and contained in $(72)$ as in the source.
Use the source's canonical prime pools, one common annular slot profile,
and individual allowed relative ray characters. A common additional
puncture of norm at most $U^B$ is allowed. It is independent of the
summed row, but may depend on frozen outer labels; the common Fourier
majorant below is taken before those labels are summed.

Write $U=Z^M$, $A=n_1+n_2+z$, and $C=2249/500$. Suppose
\[
 \frac{41M}{100}\le n_i\le\frac{11M}{25},\qquad
 \frac{3M}{100}\le z\le\frac M{25},\qquad
 n_1+n_2+Cz\le M.
 \tag{LA.1}
\]
For every requested $\epsilon>0$, sufficiently fine canonical slots
satisfy
\[
 \sum_{k\ \mathrm{eligible}}
   |S_{k,1}(Z^{n_1})S_{k,2}(Z^{n_2})Q_{k,I}|^2
 \ll Z^{M+\epsilon}\mathcal H_J.
 \tag{LA.2}
\]
The source radial majorant is understood. The factor $\mathcal H_J$
contains finite smooth seminorms and polynomial costs of the independent
plain and slot heights. The same assertion holds for the four derivative
profile pairs used by two-variable scale-supremum recovery.

The mesh can be selected from the requested loss and fixed bounded
geometric ranges, before the eventual finite slot count and live labels.
For the fixed arithmetic family and slot system, the finite profile and
height orders, constants, and eventual threshold are selected before
the live scales, masks, row selections, and heights. This is the
fixed-family quantifier order proved here. No assertion that the compiled
new density adapter already chooses its height order before every fixed
character modulus is needed or made.
\end{proposition}

\begin{proof}[Proposed source proof, under review]
One can normalize the argument to base $U$ and $M=1$; we retain $M$ in
the displayed ledgers to make their relation with the source transparent.
Put
\[
 a=\frac{M}{2000},\quad \sigma=3a,\quad k_0=\frac7{10},\qquad
 g=\bigl(n_1+n_2+\tfrac92z-M\bigr)_+.
\]
The elementary consequences of (LA.1) are
\[
 g\le\frac M{12500}<\frac M{10000},\qquad
 \frac{85M}{100}\le A\le\frac{44753M}{50000}<\frac{9M}{10}.
 \tag{LA.3}
\]
All small positive losses denoted below by $\eta$ are allocated from a
single prescribed budget. That budget is smaller than a fixed fraction
of both $\epsilon$ and $4M/225000$. The mesh is smaller than $a/10$
and a sufficiently small fraction of this budget. In particular one
greedy overshoot fits this budget. After the finite slot list has been
fixed, its aggregate support tolerance tends to zero as $Z\to\infty$;
it is not charged once per slot.

\paragraph{1. Comparison with the same masks and rectangles.}
Let $L=M/4$. Both original plain lengths exceed $L$. Form the two
equal-product rectangles with comparison lengths
$L,n_1+n_2-L$, exactly as in B3.50. The whole-product character,
the two profiles, the original slot coefficients, and every common
puncture are identical in the two rectangles. This equality is
retained through every later extraction.

To bound the comparison separately, expand only its additional common
mask by B3.41. The absolute mass of the two divisor normalizers and
the frozen prime-slot coefficients is $Z^{o(1)}$. This follows uniformly
for polynomial-size masks from
\[
 \prod_{p\mid R}(1-(Np)^{-1/2})^{-1}\ll_\eta Z^\eta.
\]
Every resulting first natural plain scale has length at most $M/4$;
clipping a nonempty subunit scale to zero does not increase that bound.
Reflect the second natural factor only if its length exceeds $M/2$.
The combined conductor bound B3.42, its exact deletion/restoration
identity B3.44, and the source scale-supremum argument put every retained
reflected length below $M/2+\xi$. The unreflected second factors are
already below $M/2$. The same eligible rows remain nonprincipal for
this reflection. Mask invariance preserves their literal exclusion;
deleting slots only restricts the permitted positive family, or leaves
the nonprincipal zero-slot family if all slots disappear.

Thus every comparison daughter has incumbent budget at most
\[
 M/4+M/2+\xi+(9/2)(M/25)=93M/100+\xi<M.
\]
Use the existing $9/2$ moment, and sum the absolute coefficient masses
and common-profile scale boxes. All independent heights remain in
the profile controls. This proves the comparison estimate uniformly.
The original centered difference still has its original whole-product
mask; no indicator is inserted into a Poisson summation formula.

\paragraph{2. First transform and the joint first-label Cauchy step.}
Enlarge the positive norm of the whole centered difference to the source
smooth row ball, as in B3.5. Its first zero frequency is bounded by
$Z^{M+\eta}$ using the powerful-product count in B3.6. Apply the exact
first transform to the remaining frequencies, including the complete
first coprimality projector and its signs. Work in a fixed dyadic sector
of the common widths $c,d,p,R,E,s_0$. Here $p$ is the first common
radical width and $s_0$ is the first projector width. Retain the common
labels, including that projector, inside the Hilbert-space direct sum
when applying Cauchy.

Let $\lambda$ denote those labels and $\mathcal N_{C,\lambda}$,
$\mathcal N_{D,\lambda}$ the exact positive Gauss norms B3.54.
Set $B_c=(3c-5d-R)_+/6$ and define $B_d$ by exchange of sides.
The outside power from the first transform is
$Z^{M-A-R/2-E}$. It suffices to prove
\[
 \sum_\lambda\mathcal N_{C,\lambda}
 \ll Z^{p+A-c+R+E+B_c+\eta},
 \qquad
 \sum_\lambda\mathcal N_{D,\lambda}
 \ll Z^{p+A-d+R+E+B_d+\eta}.
 \tag{LA.4}
\]
Indeed joint Cauchy then gives total exponent
\[
 M+p-\frac{c+d}{2}+\frac R2+\frac{B_c+B_d}{2}+\eta\le M+\eta
\]
by the primewise balance B3.56. No independent count of $\lambda$
is made after (LA.4). For the branches treated per label below, the
old allowance contains $-s_0$; the elementary first-label count
$Z^{p+s_0+\eta}$ supplies exactly (LA.4). For the high branch the
joint count will give (LA.4) directly.

The common Fourier measures are taken before the live labels, using the
ambient-box construction of B3.3 and the transport supplement. If a
continuous profile parameter varies with an outer label, use the common
integrable Fourier-density majorant provided there; do not select a
different uncontrolled Fourier measure inside each summand.

\paragraph{3. Amplification and the actual main radius.}
Use the auxiliary good primes of norm in $(Z^a/2,Z^a]$, disjoint from
the live slots. Deleting primes in the row, common supports, masks,
or first projector removes $O(\log Z)$ choices. The averaging and
Gauss identities B3.58 give one main amplified norm with density
$Z^{-a+\eta}$ and the three extracted error alternatives of B3.59.
The normalized prime average in those errors remains normalized.

For the main norm use the genuine smooth ball with exponent
\[
 K_{\rm main}=K_0+(d-c)_++6a+\omega,\qquad
 K_0=2A-c-d+R+E-M.
 \tag{LA.5}
\]
Here $\omega>0$ includes the first retained-frequency padding and
the fixed support-doubling reserve. Since the retained first dyad
obeys $K\le K_0+\delta_1$ and $\omega\ge\delta_1$, this ball
covers the amplified range of exponent $K+6a$, including its fixed
support factors for sufficiently large $Z$. More precisely this is
the existing \path{mainCommonRadius}, with $\sigma=3a$ and
$\omega=\delta+\mathrm{reserve}$; \path{main_energy_common}
proves the majorization. Thus (LA.5) changes the actual smooth ball,
not merely its declared upper bound. The Poisson identity B3.60 is
valid at this ball radius. The extracted errors retain their old
unamplified radii.

In the main second transform put $j=G_cVh'$, with the genuine
common-support partition of B3.62--B3.63. Its nominal child width is
\[
 M'=M-(c-d)_+-6a-\omega-g_2+t_2\le M-6a-\omega,
 \tag{LA.6}
\]
since $g_2\ge t_2$. For the C-side pre-allocation total
$A'=A-c-c_2$,
\[
 (A'-M')-(A-M)
 =6a+\omega-\min(c,d)-(c_2-g_2)-t_2\le6a+\omega.
 \tag{LA.7}
\]
These are direct substitutions for the actual radius; they do not
require an unchanged old unpadded radius. The nominal inner allowance
after the exact extraction normalizations is still
$M'+\Delta$, where
\[
 \Delta=b_2-p_2+B_c+a\ge a.
 \tag{LA.8}
\]

\paragraph{4. Literal children, pointwise gates, and low bins.}
Keep the full second projector. Its exact allocation B3.65--B3.69
leaves canonical slots with their original relative characters and
the same whole-product character on each side. It does not insert
new punctures on the unrestricted plain quotients. The first projector
divides both genuine second common parts; the second projector has
no such divisibility assertion.

The moving child character $\tau$ is distinct from the fixed slot
anchor $\eta_0$. Its fixed exclusion ideal remains $Q_*$, not an
ideal recomputed by inserting $\tau$ into the anchor. The source
relative-character dictionary and \path{canonical_all_twists}
show that each ray presentation actually used in the prime bounds
is nonprincipal on canonical kept rows. The plain presentation is
nonprincipal by the corresponding original row gate. Mask transport
preserves both the exclusion and every positive-real zero maximum.
Choose the finite dictionary of these presentations and the fixed
reciprocity/unit sectors before live labels. Each zero bin is formed
within one fixed parent translate $\tau_0\operatorname{Char}(Q_*)$;
different parent translates are partitioned first, at a fixed finite
cost. Conjugation of a whole rectangle is conjugation of its whole
character, whose zero maxima agree after height reflection. In
particular $\overline{\eta\psi_w}$ must not be replaced by
$\overline\eta\psi_w$. These identifications and the precise coset
scope are proved in \path{FixedDictionaryBridge.tex}.
Principal rows are handled in paragraph 7, not by a high-bin
pointwise estimate. No arbitrary additional twist is inserted into
this nonprincipal dictionary without its own gate.

Attach the common buffered bins to the actual effective row described
in paragraph 5. A low bin has upper zero parameter at most
$k_0+O(e)$. The original pointwise prime estimate, applied first to
the natural nonprincipal presentation and then transported across
the polynomial-size Euler holes, bounds removal of slots of total
width $d_z$ by squared cost
$Z^{(k_0+O(e))d_z+\eta}$. This applies to the actual canonical prime
sums, including the original nonunit zeros.

The indicator of this low bin is imposed only after both transforms,
on a positive child energy. The identities
\path{AllocatedChildren.commonChild_energy},
\path{allocated_radial_bound}, and
\path{AllocatedRayDictionary.common_clipped_energy}
allow this arbitrary additional \path{keep} predicate.
Pointwise remove one greedy prefix until the existing $9/2$ capacity
holds, or until no slots remain. Then enlarge the remaining positive
energy to the full canonical eligible family and use the incumbent
moment. In the exhausted case use its unrestricted zero-slot
estimate, which does not require the ordinary sum of the plain
lengths to be below the width. No Poisson formula is applied to the
low-bin indicator and no improved descendant moment is invoked.

There is a small but necessary width correction. Let $d_*$ be the
actual canonical width minus $M'$. Its upper bound is the declared
support padding $\delta_2$, but it can be negative if displayed
moving support cancels. If $F$ is the nominal affine defect, then
the actual defect is at most $(F-d_*)_+$. For every $0\le r\le1$,
\[
 d_*+r(F-d_*)_+\le\delta_2+rF_+.
 \tag{LA.9}
\]
This follows by considering the signs of $F-d_*$ and $F$; it is
\path{width_plus_low_defect}. Apply it with
$r=(k_0+O(e))/(9/2)$. Equations (LA.7) and B3.67--B3.69 give the
nominal defect bound $F\le g+6a+\eta$. Consequently the combined
width and slot-removal cost, including one greedy mesh increment,
is at most
\[
 M'+\frac7{45}(g+6a)+\eta.
\]
It fits (LA.8), since
\[
 a-\frac7{45}(g+6a)
 \ge\frac M{2000}-\frac7{45}
       \left(\frac M{10000}+\frac{3M}{1000}\right)
 =\frac{4M}{225000}>0.
 \tag{LA.10}
\]
Here $\omega$, $O(e)$, the support padding, and the one mesh
overshoot all belong to the already chosen smaller $\eta$ budget.
This proves the per-label low contribution within the old allowance;
its first-label count gives (LA.4).

\paragraph{5. The main high bins are counted jointly.}
Use the signed local identities from the reconstruction supplement.
For the relevant first side the actual row, including a possible
fixed unit, is
\[
 w=eF_cG_cVh',\qquad
 L_c=\log_Z NF_c,\qquad T_c=(c-d)_+.
\]
The active residues in $F_c$ are $(1+t_p)\bmod6$ on the formal
left and $(1-t_p)\bmod6$ on the formal right. A zero residue retains
its old puncture. Let $R_0$ be the canceled active radical width.
The actual frequency support gives
\[
 \log_Z Nw\le W+O(\eta),\qquad
 W=M-R-T_c+L_c-6a.
 \tag{LA.11}
\]
The favorable extra $-\omega$ can be discarded here. Every genuine
second common prime and the first projector divides $w$. The first
common primes absent from $w$ have total width at most
$p-R-E+R_0$. At fixed $w$ and their missing radical, all first and
second common labels and the first projector are divisor-many, by
the arbitrary-valuation reconstruction lemma. This is the actual
signed source identity, not a guess based on a finite valuation table.

For $\delta\in[7/10,3/4]$, put
$\rho=1-10\delta/[3(5-2\delta)]$. The proved unmarked density
has a common scale threshold before height, and the sixth-power
lifting is convergent because $\rho\ge2/7>1/6$. It therefore
counts these joint labels and rows by
\[
 Z^{\rho W+p-R-E+R_0+\eta}(1+T_h)^J.
 \tag{LA.12}
\]
A nonempty row range with slightly negative nominal $W$ has
$W\ge-O(\eta)$; apply density at the supported bound and charge
this same tolerance. If the supported bound is negative, the range
is empty. One does not use a population upper bound smaller than
one for a nonempty bounded window.

The second projector label is not included in (LA.12). Allocate it
exactly as in B3.65. Its dyadic count is paid by the two extraction
scalars B3.69. The actual pointwise plain and slot estimates give
paired normalized size at most
$Z^{\delta(A-c-b_2)+\eta}$ before these reductions; each additional
post-allocation reduction has a nonpositive contribution. Explicitly,
B3.68 gives paired pointwise exponent
$\delta(\alpha_1+\alpha_2-r_{\mathrm{clip},1}
-r_{\mathrm{clip},2})/2+O(\theta_N)$.
B3.69 already pays the divisor dyad and both extraction scalars;
the remaining $-\delta(r_{\mathrm{clip},1}+r_{\mathrm{clip},2})/2$
is nonpositive. This justifies the displayed normalized exponent
without treating the second projector as a divisor of $w$. The outside
power, from B3.60 with (LA.5), is
\[
 A-\min(c,d)+R+E-M+5a+g_2-t_2-b_2+O(\eta).
\]
Relative to the joint target $p+A-c+R+E+B_c$, the resulting exponent is
\begin{align*}
 &M(\rho-1)+\delta A+a(5-6\rho)\\
 &\quad +(1-\rho)T_c-\delta c-(1+\rho)R-E-B_c
         +\rho L_c+R_0\\
 &\quad +g_2-t_2-(1+\delta)b_2+O(\eta).
 \tag{LA.13}
\end{align*}
The second line is at most
$(4\rho-1-\delta)p-E-B_c\le-79p/270-E-B_c$:
use $c=\min(c,d)+T_c$, $\min(c,d)\ge p\ge R$,
$L_c\le5(R-R_0)$, $\rho+\delta>1$, and $5\rho>1$.
The third line without errors is at most $-\delta b_2-t_2$.
Finally (LA.3) implies the weaker bound $A/M\le113/125$, and
\[
 1-\rho-\delta A/M\ge259/16875,
 \qquad 5-6\rho\le23/7.
\]
Thus, even allowing total high-branch errors $M/1000$, the first
line and those errors are smaller than $-M/100$. Every bin straddling
$7/10$ is assigned to the low branch, using its upper endpoint
$7/10+O(e)$ in the pointwise removal cost. The high branch contains
only bins with lower endpoint at least $7/10$, so its zero-real-part
parameter is at least $17/20$, exactly within the stated density
range. The low-branch $O(e)$ rounding cost fits the reserved budget.
This proves the main high part of
(LA.4), directly with joint labels. There is no later $Z^p$ or
$Z^{s_0}$ multiplication.

\paragraph{6. The extracted amplification errors.}
Retain the old error radii and use the fine form of B3.70.
The three alternatives have
\[
 (w,w_o)=(\lambda,\lambda),(6\lambda,0),(7\lambda,\lambda),
 \qquad \lambda\ge\sigma/4
\]
for all sufficiently large $Z$. The nominal defect is at most
$g+\sigma+w-w_o+\eta$, whereas the inner allowance has
$\Delta\ge w$. Use global pointwise slot cost $3/4$ and incumbent
capacity $9/2$, so the removal rate is $1/6$. The same width
correction (LA.9) applies. The reserves are at least
\[
 13M/120000,\quad193M/120000,\quad119M/60000,
\]
respectively, before small losses. These are obtained by substituting
$g\le M/10000$, $\sigma=3M/2000$, and $\lambda\ge\sigma/4$.
The moving factors at valuations one and seven and every extraction
scalar are retained. The normalized prime average costs no second
prime count. No low-bin heredity or new zero-density statement is
used in these branches.

\paragraph{7. Zero frequencies and fixed-inducing children.}
The artificially added zero Gauss row is unamplified and has the
direct B3.57 bound, including $-s_0$. The second zero frequency
has, for (LA.5), excess
\[
 A-M-\min(c,d)-B_c+5a+O(\eta)\le A-M+5a+O(\eta)<0.
\]
For fixed-inducing nonzero second rows, retain the paired exceptional
predicate, forced congruence residues, raw extraction scalars, and
common masks. B3.73--B3.79 apply with the actual smooth radius.
For the main norm substitution into B3.75 gives exactly, before
the $\omega$ padding,
\[
 F_1=\begin{cases}
 (c+5d+R+E)/6+B_c-4a,&c\ge d,\\
 c+(R+E)/6+B_c-4a,&c\le d.
 \end{cases}
\]
The padding changes this by $-5\omega/6$. Thus
$F_1\ge2c/3-4a-O(\eta)$ and $F_2\ge2b_2/3$.
The earlier $-5a$ expression is a valid weaker lower bound, not
the exact equality: it drops the positive amplifier density term
$\ell=a$ in B3.75.

With $v=c+\min(c_2,d_2)$ and $L=M/4$, the exact maximum
\[
 \max_{v\ge0}\{A-5M/6-2v/3-(L-v)_+\}=A-M
\]
is attained at $v=L$. The main exceptional excess is consequently
at most $A-M+4a+O(\eta)<0$. The unchanged error radii use the
old bound $A-M+3\sigma+O(\eta)<0$. The mask-uniform lattice
cancellation B3.77 cancels the two leading terms because the
rectangles have the same product, character, mask, and whole-product
norm power. Independent input heights are part of the two common
profiles; they do not destroy that equality of leading terms.
No exceptional contribution is hidden in an arbitrarily small
error despite the fixed positive amplifier length.

\paragraph{8. Heights, losses, and recovery.}
The density theorem chooses its internal detector exponent first and
returns a finite height degree. Choose a distinct external bin exponent
$\tau_{\rm out}>0$ only afterward, so every occurrence of
$(1+Z^{\tau_{\rm out}})^J$ fits the reserved power losses. For
bounded input heights below a fixed small power of
$Z^{\tau_{\rm out}}$, retain the separated frequencies in a still
smaller part of that height buffer. All other common Fourier tails
are paid with arbitrary integration-by-parts orders, using raw
polynomial bounds before any arithmetic enlargement. Those orders
increase only the finite seminorm and allowed input-height powers.

Arbitrary larger input heights do not require a detector at an
uncontrolled height: absolute counting gives a fixed polynomial
bound in $Z$, and a sufficiently large fixed power of the largest
input height absorbs it. This supplies the stated polynomial height
allowance. The arithmetic scales have fixed bounded exponent ranges,
so one such power suffices. The height threshold precedes the
eventual choice of $Z$; it is not selected separately for each row.

There are only the two transforms and the calls to already established
moments. Choose smooth derivative and height orders backward through
this finite graph, then the Fourier-tail orders, and finally the
eventual threshold. Coefficient divisor bounds and the finitely many
dyadic sectors receive small shares of the original $\eta$ budget.
The first-label Cauchy step averages the two bounds (LA.4), and the
inner Cauchy step averages its two child costs; neither duplicates a
numerical loss. Minkowski adds coefficient masses already charged
in the displayed estimates. The aggregate $\eta$ losses are at most
the requested $\epsilon$.

Combining the comparison, zero-frequency, low, high, extracted-error,
and exceptional estimates proves (LA.2). Replacing either plain
profile by its logarithmic scale derivative repeats the same argument
with finite increased seminorm orders. The source's two-variable
Sobolev recovery therefore gives the four derivative-profile pairs
and the stated bounded scale supremum. This proves the proposition.
\end{proof}

\paragraph{Scope of this assembly.}
This is an analytic source proof relative to the pinned incumbent
theorems and the companion derivations explicitly named above. It is
not a claim that a new complete Lean endpoint or any independent
checker receipt exists. In particular the fixed-family quantifier
order is stated explicitly: the proof supplies common exponent savings
and a mesh chosen before live arithmetic labels, while constants and
finite height orders may depend on the fixed arithmetic dictionary.
To use this proposition in the endpoint argument, combine it with
the exact endpoint certificate, the downward-scale dichotomy,
the original weighted Euler coefficient mass, and the original
target-uniform supremum contradiction. Those are separate from the
local-moment estimate itself.

% Development supplement, 11 October 2026.
% The arithmetic dichotomy is unconditional.  The application still assumes
% the proposed local physical moment and the pinned incumbent moment.
\section{Shortened physical scales do not require a new local hypothesis}
\label{sec:physical-scale-shortening}

The local moment window imposes a lower bound on each plain length.
Consequently the sentence ``the varying conductor only shortens the
reflected scale'' does not, by itself, justify applying that moment.
The following dichotomy closes this particular gap. It uses the
\emph{existing} $9/2$ moment on the shorter scales.

Set
\[
 h=\frac{3969737987}{5000000000},\qquad
 \ell=\frac{626701}{3750000},\qquad y=h-2\ell,
 \qquad C=\frac{2249}{500}.
\]
For a row exponent $d\in[h-10^{-3},h+10^{-10}]$ and a free-shift
exponent $t\in[0,10^{-3}]$, normalized relative to $U=Z^d$, the
upper reflected plain length is
\[
 n(d,t)=1-\frac yd+t.
\]
Since $y>0$, this is increasing in both $d$ and $t$. The exact
endpoint values used in the scalar audit are
\[
 n_- =\frac{4998608000}{11894213961},\qquad
 n_+ =\frac{402041377237}{952737117000},
\]
and the nominal selected prime widths lie between
\[
 z_- =\frac{1-2n_+}{C}
      =\frac{74327181263}{2142705776133},\qquad
 z_+ =\frac{1-2n_-}{C}
      =\frac{948498980500}{26750087198289}.
\]
Direct rational comparison gives
\[
 \frac{41}{100}<n_-\le n_+<\frac{211}{500}<\frac{11}{25},
 \qquad
 \frac3{100}<z_-\le z_+<\frac{71}{2000}<\frac1{25}.
\]
All comparisons are strict, so sufficiently small fixed annular,
box, fixed-conductor, and selection-mesh losses fit inside them.
The numerical margins before these losses are approximately
$0.0000143651$ for the upper plain bound $211/500$ and
$0.0000422098$ for the upper slot bound $71/2000$.
In particular those losses must be chosen below these margins;
merely fitting inside the larger local window is insufficient for
the quantitative short-scale reserve below.

\begin{lemma}[A downward scale dichotomy]
Let $n,x_1,x_2,z\ge0$ satisfy
\[
 x_i\le n\le\frac{211}{500},\qquad
 \frac3{100}\le z\le\frac{71}{2000},\qquad
 2n+Cz\le1.
\]
Then either both actual plain lengths $x_i$ belong to
$[41/100,11/25]$ and
\[
 x_1+x_2+Cz\le1,
\]
or
\[
 x_1+x_2+\frac92z\le1-\frac{33}{4000}.
\]
\end{lemma}
\begin{proof}
If both $x_i\ge41/100$, their upper bounds follow from
$211/500<11/25$, and their new capacity follows by comparison
with $2n+Cz$.
If, for example, $x_1\le41/100$, then
\[
 x_1+x_2+\frac92z
 \le\frac{41}{100}+\frac{211}{500}
       +\frac92\frac{71}{2000}
 =1-\frac{33}{4000}.
\]
The other case is symmetric. Equality at the boundary is harmless.
\end{proof}

Here is the precise analytic application. Assume the proposed local
moment has been established for its stated profile and canonical-slot
family, and retain the pinned source's existing $9/2$ moment.
Fix a selected slot family of total width $z$ obeying the lemma,
with the same natural row zeros and eligibility in both estimates.
For any actual pair of reflected scales bounded above by $U^n$, clip
nonempty subunit scales to length zero using the source's bounded
annular dilation rule. Apply the local moment on the rectangle
$x_1,x_2\ge41/100$; on either complementary rectangle apply the
incumbent moment. Both bounds have the same target $U^{1+\epsilon}$.

The original scale-supremum recovery applies separately on these
three rectangles. In logarithmic scale coordinates it bounds the
supremum by integrals of the function and its first derivatives in
each coordinate. These are exactly the four profile pairs appearing
in the two-variable recovery. Their supports remain fixed annuli;
the number of scale boxes is polynomial in $\log U$, and is absorbed
in an arbitrarily small part of $\epsilon$. Thus row-dependent
conductor deficits, downward Euler-restoration scales, and the low
annuli of the reflected Schwartz profile do not force a local-moment
claim below $41/100$. The upper annular tail is first removed with
the original rapid-decay estimate.

The algebra relating this to the actual numerator is the pinned
source's B0.42--B0.44: write a monomial shift as $h_0h_r$, with
$h_r$ charged to the repeated-prime conductor deficit. Its reflected
scale is bounded by a fixed constant times $(U/Y)h_0$; natural Euler
restoration only shortens that scale. For a fixed $h_0$ dyad this
is the upper bound $n(d,t)$ above. The same combined conductor must
be used in this bound and in the coefficient damping. No factor of
the conductor deficit may be charged twice. The argument introduces
no further coefficient mass and therefore preserves the source's
weighted free-shift mass $\sum |c_u(h_0,h_r)|h_0^{\delta/2}$.

\paragraph{Short actual prime supply.}
There is a second possible shortening: the actual selected prime width
$z$ can be smaller than the nominal selection capacity
$z_A=(1-2n)/C$. Error slots carry zero gain, and even a main slot
may carry zero gain. No lower bound on the positive-gain supply is
required. If $z\le3/100$, the old moment applies directly, since
\[
 x_1+x_2+\frac92z
 \le 2\frac{211}{500}+\frac92\frac3{100}
 =1-\frac{21}{1000}.
\]
Thus every actual selected width $0\le z\le z_A\le71/2000$ is
covered: use this old-moment branch for $z\le3/100$, and the
preceding downward-scale dichotomy otherwise. The common worst
reserve in either old-moment branch is $33/4000$.

The selection gain is still the gain at the \emph{target capacity}
$z_A$, not merely at the smaller actual width $z$. To see this
from the literal source, write $w_i$ for every original slot width,
$L=\sum_iw_i$ for their total, and $g_i\in[0,\delta/2]$ for the
main-masked amplitude-bin gains. The error gains are zero but their
widths are retained in $L$. Set
\[
 q=\frac{\sum_iw_i g_i}{L}.
\]
The source theorem \path{zero_error_whole_slots} assumes only
$0\le z_A\le L$, $w_i\le\mathrm{mesh}$, and the above gain bounds.
It returns a subset $I$ of the positive-gain main slots satisfying
\[
 z=\sum_{i\in I}w_i\le z_A,\qquad
 \sum_{i\in I}w_i g_i\ge qz_A-\frac\delta2\mathrm{mesh},
 \qquad
 \left|\prod_{i\in I}Q_i\right|^2
 \ge U^{2qz_A-\delta\mathrm{mesh}}.
\]
There is no premise that the positive-gain main supply is at least
$z_A$. If it is exhausted first, the selected weighted gain is
$qL\ge qz_A$, which is stronger than the stated rounded bound.
Zero-gain main factors must not be reinserted into this lower-bounded
product; such factors could be arbitrarily small or zero.

The source's \path{source_class_main_selection} constructs the
amplitude classes on all original slots and uses exactly this masked
mean. Consequently the same $I$, $q$, and lower bound work uniformly
on its class of rows. Its full-main \emph{upper} bound is unchanged.
To substitute the local capacity for its old capacity one invokes
\path{zero_error_whole_slots} at $z_A$ directly, after checking
the total-original-supply inequality $z_A\le L=\ell/d$. This is the
scalar audit's capacity gate; it does not assert a fictitious
positive-main supply lower bound.

Combining this selection with either applicable physical moment
gives the same denominator gain $U^{2qz_A}$, up to the fixed mesh
loss, whether the actual width is in the local window or below it.
The old-moment branch therefore preserves the desired weighted
numerator saving. This argument supplies neither the new local
moment nor its literal initial-state adapter; it closes the
downward-scale and short-selected-supply parts conditional on those
inputs and the pinned incumbent moment.

The six elementary statements in \path{ScaleShortening.lean}
check the rational upper edges, both strict incumbent reserves,
the two dichotomies, and the exhausted-supply affine gain. Their
scope is arithmetic; the actual selection and spike identity are
the source theorems identified above.

\paragraph{All larger free shifts, including zero capacity.}
For completeness the restriction $t\le10^{-3}$ does not discard the
remaining coefficient dyads. Put $b=2y/d-1>0$, $t_0=10^{-3}$,
and
\[
 z_{\mathrm{new}}(0)=\frac bC,\qquad
 z_{\mathrm{old}}(t)=\max\left\{0,\frac{b-2t}{9/2}\right\}.
\]
The actual incumbent capacity uses $6\kappa\le9/2$, so this is a
lower admissible capacity. The weighted coefficient mass in B0.44
pays $U^{-\delta t/2}$. Thus it suffices to prove
\[
 \frac{\delta t}{2}+\frac{qz_{\mathrm{old}}(t)}2
 \ge\frac{qz_{\mathrm{new}}(0)}2.
 \tag{SS}
\]
Write
\[
 S_b(t)=\frac t2-\frac14\bigl(b/C-z_{\mathrm{old}}(t)\bigr).
\]
This function is continuous at $t=b/2$. Its derivative is $7/18$
where the old capacity is positive and $1/2$ where it is zero.
For the audited near-endpoint rows,
$b\le b_{\max}=1-2n_-$ and $b>2t_0$.
The exact rational audit gives
\[
 S_b(t_0)\ge
 \frac{t_0}{2}-\frac14
   \left(\frac{b_{\max}}C-
         \frac{b_{\max}-2t_0}{9/2}\right)
 >\frac3{10000}.
\]
It follows that $S_b(t)>0$ for every $t\ge t_0$, with no exception
at the clipping point. Since $0\le q\le\delta/2$ and
$z_{\mathrm{new}}(0)-z_{\mathrm{old}}(t)\ge0$, the left side of
(SS) minus its right side is at least $\delta S_b(t)\ge0$.
This proves the full large-shift payment, including the empty selected
family at zero capacity. That family uses the incumbent unrestricted
zero-slot moment B3.40, not a moment with a hidden requirement
$x_1+x_2\le1$. In the source definition \path{ZeroAt}, the
positive scales are required only to satisfy $X_i\le Z^L$; its
conclusion is $Z^{s.\mathrm{width}+\epsilon}$. The actual physical
specialization \path{finite_pair_zeroAt} in
\path{Cycle25/Energy/Physical/CertifiedZero.lean} gives
$U^{1+\delta+\epsilon}$ for $X_i\le U^2$.
Its \path{certified_reflected_zero_fourier_moment}
retains both independent Fourier heights.
Finally \path{certified_original_numerator_fourth} in
\path{Cycle25/Energy/Physical/OriginalFourth.lean} accepts
the exact disjunction
\[
 2(m+\xi)+6\kappa\sum_{i\in I}w_i\le1
 \quad\text{or}\quad I=\varnothing,
 \qquad m+\xi\le2,
\]
and gives the same $U^{1+o(1)}$ fourth-moment bound in either case.
There is consequently no extra positive-part length cost when the
capacity vanishes. The larger $\max(1,2m)$ bound in the unmarked
density adapter allows its finitely many exceptional rows; it is
not the eligible physical zero-slot estimate being used here.

For $0\le t\le t_0$ the local capacity is
$z_{\mathrm{new}}(t)=(b-2t)/C>0$.
Here the same conclusion follows directly from
\[
 \frac{\delta t}{2}+
 \frac q2\bigl(z_{\mathrm{new}}(t)-z_{\mathrm{new}}(0)\bigr)
 =t\left(\frac\delta2-\frac qC\right)\ge0.
\]
The original support cutoff restricts $t\le m+o(1)$, as in B0.44;
all scales therefore stay in bounded geometric ranges. The new
argument requires no new uniform estimate for unbounded plain
lengths. Row exponents below the near-endpoint interval use the
separate incumbent row-endpoint certificates, as specified in the
global endpoint assembly.

% Proposed complete written endpoint assembly, 11 October 2026.
% This source proof is under review. It is not a new Lean endpoint theorem,
% a clean replay, or an independent-checker certificate.
\section{The endpoint and the target-uniform continuation argument}
\label{sec:endpoint-assembly}

We now insert the local physical moment into the pinned numerator
construction. The input results are the established incumbent moments,
the exact physical Euler expansion and coefficient mass in B0.32--B0.44
and B4 of the pinned Cycle25 manuscript, its detector count B0.40, and
its contour and continuation construction B0.7 and T2--T3. The new input
is the local-moment proposition proved in
\path{PhysicalLocalAssembly.tex}, with the fixed-dictionary and
height-uniform density arguments preceding it. The downward-scale
argument is \path{PhysicalScaleShortening.tex}. We retain the exact
rational certificate in \path{local_moment_endpoint_audit.py}.
No unproved varying-parent moment is introduced in this application.

Set
\[
 \theta=\frac{4374433}{5000000},\quad
 \ell=\frac{626701}{3750000},\quad
 h=\frac{3969737987}{5000000000},\quad
 x=1+\ell-h,\quad y=h-2\ell,
 \qquad C=\frac{2249}{500}.
\]
In particular $\theta=11/12-\ell/4$, $x+y=1-\ell$,
$0<\ell<x<y<2x$, and $y<1/2<h<2y$.

\begin{theorem}[Proposed endpoint, relative to the pinned analytic framework]
Every finite-order Hecke character of the Eisenstein field with trivial
infinite type has no zero in $\Re s>\theta$, with the principal pole at
$s=1$ excluded. Presented characters may be imprimitive and have
arbitrary conductor. Every ordinary Dirichlet $L$-function has the same
zero-free half-plane, again excluding the principal pole. In particular
$\zeta(s)\ne0$ there for $s\ne1$.
\end{theorem}

\begin{proof}[Proposed source proof, under review]
Let $\beta_*$ be the source's Hecke zero supremum, including its sentinel
$1/2$. The incumbent theorem gives $\beta_*\le7/8$. Suppose
$\mu=\beta_*-\theta>0$ and set
\[
 \kappa=2\beta_*-1,\qquad
 \kappa_*=2\theta-1=\frac{1874433}{2500000}.
\]
Then $\kappa_*<\kappa\le3/4$. The incumbent marked detector and ordinary
moment apply at this $\kappa$: their hypothesis $2\beta_*-1\le\kappa$
holds by equality. The new local moment uses only the incumbent
$9/2$ and unrestricted zero-slot estimates at descendant states.

\paragraph{1. The actual physical family and selected slots.}
Fix a primitive target $\eta$ only after the common real margins and slot
mesh have been chosen. Its finite exclusion set $S$, original local
Euler factors, and all natural row zeros are those of the source probe.
In the physical numerator moment the fixed base is
\texttt{fixedSourcePrincipal $S$}; the variable part is the sextic row
character. This is exactly the coefficient identity in
\path{PhysicalMomentInput.source_physical_slot_conj} and
\path{reflected_pair_norm_eq}. The relative prime coefficient
contains the inverse fixed base and the common conjugate annular
profile, with real power $33/50$ and individual height $-\Im z_i$.
Conjugation applies to the entire reflected product. Thus the local
proposition applies to this fixed family with the same literal
eligibility and masks. The fixed exceptional free rows have bounded
norm and are absent from every growing retained dyad; the source
\path{weighted_source_exceptional_norm} proves this fact.

In an amplitude class let $w_i=\ell_i/d$, let
$L=\sum_iw_i=\ell/d$, and let $0\le g_i\le\delta/2$ be its amplitude
exponents. Error slots have $g_i=0$. Put
$q=L^{-1}\sum_iw_ig_i=\delta p$, where $0\le p\le1/2$.
For any target capacity $0\le z_A\le L$, the exact theorem
\path{SlotSelection.zero_error_whole_slots} selects positive-gain
main slots $I$ such that
\[
 z_I=\sum_{i\in I}w_i\le z_A,
 \qquad
 \left|\prod_{i\in I}Q_i\right|^2
 \ge U^{2qz_A-\delta\,\mathrm{mesh}-o(1)}.       \tag{EA.1}
\]
The gain is measured at the target capacity $z_A$, even if $z_I$ is
strictly smaller. It is not legitimate to assume $z_I$ nearly fills
$z_A$ or to insert error slots into the selected product. The selection
is made once for the amplitude class, not independently for its rows.

Put $d_0=1/200$ and $\zeta_0=10^{-10}$. All original slots can be fixed
so short that their relative widths meet every required mesh for
$d\ge d_0$. Detector supply is available because the exact geometry
satisfies $\ell/(h+\zeta_0)>4/21$.
For $d<1/2$, the same selection theorem applies directly, without the
$d\ge1/2$ assumption of its convenience wrapper:
$L\ge2\ell>1/3$ and the old physical capacity is at most
$1/(6\kappa)\le5/21$. Thus no lower-row supply condition is missing.

\paragraph{2. The reflected scale and every free Euler shift.}
The exact coefficient expansion writes a shift as $h_0h_r$, with
$h_r$ charged to the repeated-prime conductor deficit. B0.42--B0.44 give
the actual reflected scale bound
\[
 X_u\ll_S (U/Y)h_0,
 \qquad \sum_{h_0,h_r}|c_u(h_0,h_r)|h_0^{\delta/2}
       \ll U^\varepsilon.                            \tag{EA.2}
\]
The same combined conductor occurs in both statements. Its deficit is
not charged a second time. Natural Euler restoration only shortens
scales; its coefficients have the source's bounded absolute mass.
The original support retains $h_0\ll Y$, so all scale exponents lie in
a fixed bounded range.

Consider first $h-10^{-3}\le d\le h+\zeta_0$ and
$t=\log_Uh_0\in[0,10^{-3}]$. The upper reflected length is
$n(d,t)=1-y/d+t$. Choose the target capacity just below
\[
 z_A(t)=\frac{1-2n(d,t)}{C},
\]
reserving the small capacity decrement for annular and scale-box errors.
The exact geometric checks, with strict room for those errors, give
\[
 \frac{41}{100}<n(d,t)\le\frac{211}{500},\qquad
 \frac3{100}<z_A(t)\le\frac{71}{2000},\qquad z_A(t)<L.
\]
Let $x_1,x_2\le n(d,t)$ be the actual nonnegative plain lengths, and
let $z_I$ be the actual selected width from (EA.1). There are three cases.
If $z_I<3/100$, the incumbent capacity holds with strict reserve:
\[
 x_1+x_2+\frac92z_I
 \le2\frac{211}{500}+\frac92\frac3{100}
 =\frac{979}{1000}<1.                                \tag{EA.3}
\]
If $z_I\ge3/100$ and some $x_j<41/100$, the downward-scale lemma gives
the incumbent reserve $33/4000$. In the remaining case the local
proposition applies because
$x_1+x_2+Cz_I\le2n(d,t)+Cz_A(t)\le1$.
The empty selected product uses the unrestricted zero-slot estimate.
The same two-variable scale-supremum recovery and its four derivative
profiles apply on these separate regions. Therefore the fourth moment
needed for numerator H\"older has the same $U^{1+\varepsilon}$ bound
in all cases, while its gain remains $qz_A(t)/2$ by (EA.1).

For these small shifts,
\[
 -\frac{\delta t}{2}-\frac{qz_A(t)}2
 \le-\frac{qz_A(0)}2,
\]
since $q\le\delta/2$ and $C>1$. For larger shifts use the incumbent
capacity instead. Write $b=2y/d-1$ and use the smaller comparison
capacity $z_o(t)=\max(0,(b-2t)/(9/2))$. The exact scalar audit shows
\[
 S(10^{-3})>\frac3{10000},\qquad
 S(t)=\frac t2-\frac{b/C-z_o(t)}4.
\]
While $z_o(t)>0$, its derivative is $7/18$; after the capacity clips
to zero the derivative is $1/2$. Hence $S(t)\ge0$ for every
$t\ge10^{-3}$. With $q\le\delta/2$ this proves
\[
 -\frac{\delta t}{2}-\frac{qz_o(t)}2
 \le-\frac{qz_A(0)}2.                                \tag{EA.4}
\]
The actual old capacity at coefficient $6\kappa\le9/2$ is no smaller.
Thus (EA.2) pays every free-shift dyad, including those with zero old
capacity. No Euler term has been omitted from the numerator inversion.
At zero capacity the two reflected plain lengths can have sum greater
than one; this is exactly where the \emph{unrestricted} physical MP0 is
needed. Its actual source adapter
\path{Physical/CertifiedZero.finite_pair_zeroAt} has no affine
plain-length condition and bounds the moment by $U^{1+\varepsilon}$
for each plain scale up to $U^2$. The theorem
\path{Physical/OriginalFourth.certified_original_numerator_fourth}
accordingly accepts either its affine capacity inequality or an empty
selected family. The original support bound in (EA.2) gives reflected
length at most $1+o(1)<2$, so this is a literal application of that
theorem. The weaker all-free-row plain estimate used in the new density
argument is not being substituted for this physical MP0.
Below $d=h-10^{-3}$ the original B0.41 argument is unchanged.

\paragraph{3. The common retained-row saving.}
The detector remains
\[
 R_\kappa(\delta,p)=1-\delta+
 \frac{\delta c_\kappa(p)(5/6-\delta)}
      {2(5/6-\delta+\delta c_\kappa(p))},\qquad
 c_\kappa(p)=\frac{2(1-p)(1-p/(3\kappa))}
                   {3-p-2p/(3\kappa)}.
\]
All $\delta\in[0,3/4]$ and $p\in[0,1/2]$ are included. The physical
H\"older argument and the weighted mass (EA.2), on the original full
$w$-inversion line, give the exponent $G_{C,\kappa}$ of the paper on
the local row interval. On old rows its general form is
\[
 G_o(d)=\frac{1+\delta}{2}+h\frac{13}{75}
 -\frac y2-\frac\ell2+\ell q
 +d\left(\frac14+\frac{3R_\kappa}4-\frac{17}{50}\right)
 -\frac{q}{12\kappa}\min(d,2y-d).                    \tag{EA.5}
\]
This is the exact B0.46/B0.55 ledger, retaining the reciprocal only
after the full physical numerator inversion.

At $\kappa=\kappa_*$ the rational Bernstein certificate proves savings
$10^{-7}$ at $(d,c)=(h,C)$, $10^{-5}$ at
$(h-10^{-3},6\kappa_*)$, and $10^{-3}$ at
$(1/2,6\kappa_*)$, uniformly on the full $(\delta,p)$ rectangle.
At the junction the new exponent is no larger than the old one because
$C<6\kappa_*$ and $2y-d>0$. Affinity in $d$ covers the intervening
local interval. The extension above $h$ costs at most $2\zeta_0$, so
the retained saving there is at least $499/5000000000>9\cdot10^{-8}$.

The old function (EA.5) is convex in $d$. To cover $d_0\le d\le1/2$,
combine the certificate at $1/2$ with the elementary bound at $d_0$,
using $a\le7/8$, $R_\kappa\le1$ and $q\le3/8$:
\[
 G_o(d_0)\le\frac78-\frac{49h}{150}+\frac{7\ell}{8}
                  +\frac{33d_0}{50}
 =\frac{573880513637}{750000000000}<\theta-\frac1{10}.
\]
On $[1/2,h-10^{-3}]$, affinity and its two certified endpoints suffice.
The floor classes use the source bound
\[
 \frac{51}{100}+\frac h3+\frac{51\ell}{100}
 =\frac{12898208027}{15000000000}<\theta-\frac1{100},
\]
with room for the tiny row extension. They do not require a positive
detector zero above the sentinel.

The variation in $\kappa$ must be checked separately for the two
physical capacities. On the local branch only $R_\kappa$ varies, and
the source derivatives $\partial_\kappa c_\kappa\le1/18$ and
$\partial_cR\le5/12$ give $\partial_\kappa G_{C,\kappa}<1/50$.
Because $\kappa-\kappa_*=2\mu$, its increase is less than
$\mu/25$. Thus its saving below $\beta_*$ is at least its certificate
saving plus $24\mu/25$.
For the old branch the reciprocal capacity also varies. Since
$\min(d,2y-d)\le1/2$ and $d\le1$, one has
\[
 \partial_\kappa G_o
 \le\frac5{288}+\frac{(3/8)(1/2)}{12(7/10)^2}
 =\frac{695}{14112}<\frac18.
\]
Its feedback saving is therefore at least its certificate saving plus
$3\mu/4$. In particular all retained classes have a common strictly
positive power reserve before a target character is fixed.

\paragraph{4. Principal, very small, very large, and off-height terms.}
The new geometry leaves the source contour boxes inside their stated
domains; in particular $\theta>87/100$ and the physical scalar boundary
is still exactly $11/12-\ell/4$. The principal double residue B0.35--B0.36
is the same function $f_\eta$. Its three remainders have relative
reserves $y/4$, $h/24$, and a fixed positive prime-product exponent
$\mu_P$. Their derivation uses the local Euler identity and the
principal correction, not the new moment.

For nonprincipal rows below $Z^{d_0}$ the original global-line bound
B0.56 gives, before small losses,
\[
 h\frac{13}{75}-\frac y2+\frac{63d_0}{50}
 =-\frac{21483987121}{250000000000}<-\frac1{20}.
\]
The finite nonprincipal unit rows are included. For rows above
$Z^{h+\zeta_0}$ the raw absolute estimate B0.57 applies: after choosing
$\zeta_0>0$, choose its fixed positive $z$-line parameter $v$ large
enough that the term has any prescribed negative power. This is not
an application of a central-row density estimate.

The compensated principal probe satisfies the actual B0.27 geometry:
with $M=x+y=1-\ell$ one has
\[
 y-2x+\ell=-\frac{597590039}{5000000000}<0,\qquad
 1+3\ell-2M=-\frac{123299}{750000}<0.
\]
Thus its exponent is $(5x+y)/12=\theta-k$, where
$k=5/6-x/3-\ell/6$. Division by $c_SA(Z)$ costs only a power of
$\log Z$, giving $J_\eta(Z)\ll_\eta Z^{\theta-k+\varepsilon}$.
The three independent contour frequencies and their
joining segments are those of B0.7. Raw polynomial bounds are taken
before any separation in a tail; they give $Z^{B'}T^{-N}$ with $B'$
fixed before the tail order $N$. No contour crosses an unaccounted
nonprincipal numerator pole or a reciprocal zero.

\paragraph{5. Order of choices and continuation.}
Choose from $\mu$, the geometric and certified margins a common moment
loss, capacity decrement, amplitude mesh, and a finite slot system
meeting both the old and new moment meshes on $d\ge d_0$. Its smallest
slot fixes $\mu_P>0$. Choose common real buffers smaller than these
reserves. These choices, including the power saving below, precede
the target $\eta$.

For each target now fix its finite set $S$, dictionary, and the finitely
many internal moment, profile and Sobolev orders. They may depend on the
target. Next choose a target-dependent positive analysis-height exponent
$\tau$ small enough that all factors $(1+Z^\tau)^J$ use less than the
same reserved real margin. The internal detector exponent in the density
proof was already fixed; it is not this later $\tau$. Choose the external
tail order after $\tau$, and finally the target-dependent asymptotic
threshold. The resulting constants and thresholds may depend on the
target, but one obtains a single $\Delta>0$, independent of it, with
\[
 J_\eta(Z)-f_\eta(Z)
       \ll_\eta Z^{\beta_*-k-\Delta}.                \tag{EA.6}
\]
Here $\Delta$ is smaller than fixed fractions of $\mu$, the retained
saving, $y/4,h/24,\mu_P$, and the small-row saving. Make it also smaller
than $\mu/4$. The principal probe bound then gives
$f_\eta(Z)\ll_\eta Z^{\beta_*-k-\Delta}$.

At $Z\downarrow0$ the defining Mellin line can be moved arbitrarily far
right, as in B0.60. Therefore
\[
 \int_0^\infty f_\eta(Z)Z^{-s+k}\,\frac{dZ}{Z}
\]
is holomorphic for $\Re s>\beta_*-\Delta$. On $\Re s>1$ its exact
value is
\[
 e^{(s-5/6)^2}\frac{H_\eta(s)}{L^S(s,\eta)}.
\]
The source cutoff ensures $|H_\eta(s)-1|\le1/2$ for
$\Re s>\theta>87/100$, and the Gaussian never vanishes. Since
$\beta_*-\Delta>\theta$, holomorphic continuation excludes every zero
of $L^S(s,\eta)$ to the right of $\beta_*-\Delta$, away from the
principal pole. Restoring fixed Euler factors does not add zeros in
this positive half-plane.

This is precisely the quantifier pattern of the proved, variable-boundary
\path{Cycle25.PrimitiveProbeContract} and
\path{UniformCommonProbe} in
\path{Cycle25/Assembly/Continuation.lean}. In that file
\path{beta_le_of_common_probe} first chooses a primitive zero
sufficiently near the supremum and then invokes the target's individual
probe. Its Big-O constants and threshold occur after the target
quantifier. Thus no uniform numerical bound on growing target moduli is
needed. A target-dependent power saving would be insufficient; (EA.6)
has explicitly avoided that error. The common positive improvement
contradicts the definition of $\beta_*$. Hence $\beta_*\le\theta$.

\paragraph{6. Primitive reduction and Dirichlet transfer.}
For an imprimitive Hecke character induced by $\eta^*$, the exact
identity T2.1 expresses its $L$-function as $L(s,\eta^*)$ times finitely
many factors $1-\eta^*(\mathfrak p)(N\mathfrak p)^{-s}$.
They are nonzero for $\Re s>0$. This proves the Hecke assertion for
every presented finite-order character of trivial infinite type,
with the same pole exception.

For an ordinary Dirichlet character $\chi$, use its norm lift to the
Eisenstein field. The pinned theorem
\path{HeckeDirichlet.LFunction_eq_dirichlet_product} gives
\[
 L_K(s,\eta_\chi)=L(s,\chi)L(s,\chi\chi_{-3})
\]
away from $s=0,1$, with the original imprimitive and ramified Euler
factors included. For $\theta<\Re s<1$, both factors have finite
values and their product is nonzero by the Hecke result, so
$L(s,\chi)\ne0$. On $\Re s\ge1$ the classical ordinary Dirichlet
nonvanishing theorem applies with its principal-pole exception. The
trivial character of modulus one gives the stated zeta assertion.
These are exactly the transfers in
\path{Cycle25.dirichlet_of_hecke} and
\path{zeta_of_hecke}; they require neither primitivity nor reality
of $\chi$.
\end{proof}

\paragraph{Verification status.}
This is a complete proposed written source-level assembly, under review,
relative to the pinned incumbent analytic framework and the new local
proof. The strict rational comparisons and Bernstein certificate are
exact. Selected supporting Lean modules compile, but this file does not
create a new all-family Lean theorem, a clean dependency replay, a
canonical-target comparison, or an independent kernel certificate.


\begin{thebibliography}{9}
\raggedright
\bibitem{collet}
Hailey Collet, \emph{The quartic QRH boundary}, 10 October 2026.
Cycle25 manuscript and source, revision
\code{58344dfdbe756cf2f743da1908ffe0582418cf5b}.
\href{https://github.com/HaileyStorm/QuasiRiemannTracker/blob/58344dfdbe756cf2f743da1908ffe0582418cf5b/public/proofs/cycle25-quartic-20261010/paper.tex}{Pinned manuscript}.
Equation references B0 and B3 in the companion notes refer to this revision.

\bibitem{openai}
OpenAI, \emph{OpenAI/math}, formal mathematical development,
revision \code{fd4aeeb2ee4fc729c18d98444fed42fd0529eeeb}.
\href{https://github.com/openai/math/tree/fd4aeeb2ee4fc729c18d98444fed42fd0529eeeb}{Pinned source repository}.

\bibitem{levy}
Akash Levy, \emph{Weighted-numerator proof}, source package,
revision \code{2fd60c0926b66ea18d7436f5ed55250fd006ab5d}.
\href{https://github.com/akashlevy/QuasiRiemannTracker/tree/2fd60c0926b66ea18d7436f5ed55250fd006ab5d/proofs/akashlevy-20261009-weighted-numerator}{Pinned proof package}.

\bibitem{gao}
Zhongpai Gao, with AI assistance,
\emph{Local physical moment research candidate: 0.8748866},
11 October 2026. Source-review snapshot,
revision \code{7f9166e2d626586cc3c89853497c0e0e22e856ee}.
\href{https://github.com/Gaozhongpai/QuasiRiemannTracker/tree/7f9166e2d626586cc3c89853497c0e0e22e856ee/proofs/gao-local-moment-20261011}{Research notes, certificates, and Lean components};
\href{https://github.com/GTimG/QuasiRiemannTracker/pull/15}{draft pull request 15}.

\bibitem{gaov2}
Zhongpai Gao, with AI assistance,
\emph{Joint-label reconstruction and actual unmarked inputs},
11 October 2026, second research revision.
\href{https://github.com/Gaozhongpai/QuasiRiemannTracker/tree/gao-local-moment-source-review-20261011/proofs/gao-local-moment-20261011-v2}{Companion note and nine compiled components}.
That historical revision kept the physical factorization as an
application obligation; Draft 3 supplies the source-level arguments
in the present supplements.
\end{thebibliography}
\end{document}
